% Options for packages loaded elsewhere
\PassOptionsToPackage{unicode}{hyperref}
\PassOptionsToPackage{hyphens}{url}
\documentclass[
]{article}
\usepackage{xcolor}
\usepackage[margin=1in]{geometry}
\usepackage{amsmath,amssymb}
\setcounter{secnumdepth}{-\maxdimen} % remove section numbering
\usepackage{iftex}
\ifPDFTeX
  \usepackage[T1]{fontenc}
  \usepackage[utf8]{inputenc}
  \usepackage{textcomp} % provide euro and other symbols
\else % if luatex or xetex
  \usepackage{unicode-math} % this also loads fontspec
  \defaultfontfeatures{Scale=MatchLowercase}
  \defaultfontfeatures[\rmfamily]{Ligatures=TeX,Scale=1}
\fi
\usepackage{lmodern}
\ifPDFTeX\else
  % xetex/luatex font selection
\fi
% Use upquote if available, for straight quotes in verbatim environments
\IfFileExists{upquote.sty}{\usepackage{upquote}}{}
\IfFileExists{microtype.sty}{% use microtype if available
  \usepackage[]{microtype}
  \UseMicrotypeSet[protrusion]{basicmath} % disable protrusion for tt fonts
}{}
\makeatletter
\@ifundefined{KOMAClassName}{% if non-KOMA class
  \IfFileExists{parskip.sty}{%
    \usepackage{parskip}
  }{% else
    \setlength{\parindent}{0pt}
    \setlength{\parskip}{6pt plus 2pt minus 1pt}}
}{% if KOMA class
  \KOMAoptions{parskip=half}}
\makeatother
\usepackage{color}
\usepackage{fancyvrb}
\newcommand{\VerbBar}{|}
\newcommand{\VERB}{\Verb[commandchars=\\\{\}]}
\DefineVerbatimEnvironment{Highlighting}{Verbatim}{commandchars=\\\{\}}
% Add ',fontsize=\small' for more characters per line
\usepackage{framed}
\definecolor{shadecolor}{RGB}{248,248,248}
\newenvironment{Shaded}{\begin{snugshade}}{\end{snugshade}}
\newcommand{\AlertTok}[1]{\textcolor[rgb]{0.94,0.16,0.16}{#1}}
\newcommand{\AnnotationTok}[1]{\textcolor[rgb]{0.56,0.35,0.01}{\textbf{\textit{#1}}}}
\newcommand{\AttributeTok}[1]{\textcolor[rgb]{0.13,0.29,0.53}{#1}}
\newcommand{\BaseNTok}[1]{\textcolor[rgb]{0.00,0.00,0.81}{#1}}
\newcommand{\BuiltInTok}[1]{#1}
\newcommand{\CharTok}[1]{\textcolor[rgb]{0.31,0.60,0.02}{#1}}
\newcommand{\CommentTok}[1]{\textcolor[rgb]{0.56,0.35,0.01}{\textit{#1}}}
\newcommand{\CommentVarTok}[1]{\textcolor[rgb]{0.56,0.35,0.01}{\textbf{\textit{#1}}}}
\newcommand{\ConstantTok}[1]{\textcolor[rgb]{0.56,0.35,0.01}{#1}}
\newcommand{\ControlFlowTok}[1]{\textcolor[rgb]{0.13,0.29,0.53}{\textbf{#1}}}
\newcommand{\DataTypeTok}[1]{\textcolor[rgb]{0.13,0.29,0.53}{#1}}
\newcommand{\DecValTok}[1]{\textcolor[rgb]{0.00,0.00,0.81}{#1}}
\newcommand{\DocumentationTok}[1]{\textcolor[rgb]{0.56,0.35,0.01}{\textbf{\textit{#1}}}}
\newcommand{\ErrorTok}[1]{\textcolor[rgb]{0.64,0.00,0.00}{\textbf{#1}}}
\newcommand{\ExtensionTok}[1]{#1}
\newcommand{\FloatTok}[1]{\textcolor[rgb]{0.00,0.00,0.81}{#1}}
\newcommand{\FunctionTok}[1]{\textcolor[rgb]{0.13,0.29,0.53}{\textbf{#1}}}
\newcommand{\ImportTok}[1]{#1}
\newcommand{\InformationTok}[1]{\textcolor[rgb]{0.56,0.35,0.01}{\textbf{\textit{#1}}}}
\newcommand{\KeywordTok}[1]{\textcolor[rgb]{0.13,0.29,0.53}{\textbf{#1}}}
\newcommand{\NormalTok}[1]{#1}
\newcommand{\OperatorTok}[1]{\textcolor[rgb]{0.81,0.36,0.00}{\textbf{#1}}}
\newcommand{\OtherTok}[1]{\textcolor[rgb]{0.56,0.35,0.01}{#1}}
\newcommand{\PreprocessorTok}[1]{\textcolor[rgb]{0.56,0.35,0.01}{\textit{#1}}}
\newcommand{\RegionMarkerTok}[1]{#1}
\newcommand{\SpecialCharTok}[1]{\textcolor[rgb]{0.81,0.36,0.00}{\textbf{#1}}}
\newcommand{\SpecialStringTok}[1]{\textcolor[rgb]{0.31,0.60,0.02}{#1}}
\newcommand{\StringTok}[1]{\textcolor[rgb]{0.31,0.60,0.02}{#1}}
\newcommand{\VariableTok}[1]{\textcolor[rgb]{0.00,0.00,0.00}{#1}}
\newcommand{\VerbatimStringTok}[1]{\textcolor[rgb]{0.31,0.60,0.02}{#1}}
\newcommand{\WarningTok}[1]{\textcolor[rgb]{0.56,0.35,0.01}{\textbf{\textit{#1}}}}
\usepackage{longtable,booktabs,array}
\usepackage{calc} % for calculating minipage widths
% Correct order of tables after \paragraph or \subparagraph
\usepackage{etoolbox}
\makeatletter
\patchcmd\longtable{\par}{\if@noskipsec\mbox{}\fi\par}{}{}
\makeatother
% Allow footnotes in longtable head/foot
\IfFileExists{footnotehyper.sty}{\usepackage{footnotehyper}}{\usepackage{footnote}}
\makesavenoteenv{longtable}
\usepackage{graphicx}
\makeatletter
\newsavebox\pandoc@box
\newcommand*\pandocbounded[1]{% scales image to fit in text height/width
  \sbox\pandoc@box{#1}%
  \Gscale@div\@tempa{\textheight}{\dimexpr\ht\pandoc@box+\dp\pandoc@box\relax}%
  \Gscale@div\@tempb{\linewidth}{\wd\pandoc@box}%
  \ifdim\@tempb\p@<\@tempa\p@\let\@tempa\@tempb\fi% select the smaller of both
  \ifdim\@tempa\p@<\p@\scalebox{\@tempa}{\usebox\pandoc@box}%
  \else\usebox{\pandoc@box}%
  \fi%
}
% Set default figure placement to htbp
\def\fps@figure{htbp}
\makeatother
\setlength{\emergencystretch}{3em} % prevent overfull lines
\providecommand{\tightlist}{%
  \setlength{\itemsep}{0pt}\setlength{\parskip}{0pt}}
\usepackage{bookmark}
\IfFileExists{xurl.sty}{\usepackage{xurl}}{} % add URL line breaks if available
\urlstyle{same}
\hypersetup{
  pdftitle={social\_mobility\_data\_eda},
  hidelinks,
  pdfcreator={LaTeX via pandoc}}

\title{social\_mobility\_data\_eda}
\author{}
\date{\vspace{-2.5em}}

\begin{document}
\maketitle

\section{Social Class Mobility Data
Exploration}\label{social-class-mobility-data-exploration}

The goal of this Markdown file is to investigate the data I am going to
use as a real data illustration for the methods I am developing in my
master's thesis. My master's thesis is focused on developing mediation
analysis methods with multiple mediators when the causal ordering of the
mediators is known. These methods can be used to investigate what the
causal impact of socially driven disparities in variables like education
can have on differential outcomes for these same social groups.

The data I will be using in my thesis is based on the dataset used by
Kuha et al (2021), which uses mediation analysis methods to characterize
the extent to which differences in educational outcomes across social
classes in Britain mediate relative social class mobility, a question of
interest in the sociology literature. This data comes from three
separate samples of British citizens born in 1946, 1958, and 1970, which
gives information about a unit's social class various points in their
life, their level of education over time, their occupation over time,
their results on an intelligence test they took at age 11, and their
father's social class, among other helpful covariates. I refer to these
three separate datasets as the ``work history data''.

Kuha et al (2021) does not use the intelligence test results and instead
investigates the role of only education in mediating father-child social
class mobility. To this end, they construct a cleaned dataset that
excludes the intelligence test results and looks at their social class
and education as of middle age (age 37-38) only. They also re-weight the
survey data to account for sampele selection effects and use multiple
imputation methods to fill in missing values in the data. I refer to
this cleaned dataset as the ``social mobility data'' below.

\subsection{Setup}\label{setup}

\begin{Shaded}
\begin{Highlighting}[]
\CommentTok{\# Loading required packages}
\FunctionTok{library}\NormalTok{(dplyr)}
\end{Highlighting}
\end{Shaded}

\begin{verbatim}
## 
## Attaching package: 'dplyr'
\end{verbatim}

\begin{verbatim}
## The following objects are masked from 'package:stats':
## 
##     filter, lag
\end{verbatim}

\begin{verbatim}
## The following objects are masked from 'package:base':
## 
##     intersect, setdiff, setequal, union
\end{verbatim}

\begin{Shaded}
\begin{Highlighting}[]
\FunctionTok{library}\NormalTok{(ggplot2)}
\FunctionTok{library}\NormalTok{(haven)}
\FunctionTok{library}\NormalTok{(tidyr)}
\FunctionTok{library}\NormalTok{(skimr)}
\FunctionTok{library}\NormalTok{(GGally)}
\FunctionTok{library}\NormalTok{(janitor)}
\end{Highlighting}
\end{Shaded}

\begin{verbatim}
## 
## Attaching package: 'janitor'
\end{verbatim}

\begin{verbatim}
## The following objects are masked from 'package:stats':
## 
##     chisq.test, fisher.test
\end{verbatim}

\begin{Shaded}
\begin{Highlighting}[]
\CommentTok{\# Setting working directory}
\FunctionTok{setwd}\NormalTok{(}\StringTok{"/Users/19175/Documents/personal/research/MSc\_Thesis{-}Mediation\_Analysis/data"}\NormalTok{)}

\CommentTok{\# Loading input data}
\NormalTok{path }\OtherTok{\textless{}{-}} \StringTok{"input/social\_class\_mobility\_data/"}
\NormalTok{social\_mobility\_data }\OtherTok{\textless{}{-}} \FunctionTok{readRDS}\NormalTok{(}\FunctionTok{paste0}\NormalTok{(path, }\StringTok{"social\_mobility\_data.rds"}\NormalTok{))}
\NormalTok{workhistory1 }\OtherTok{\textless{}{-}} \FunctionTok{read\_dta}\NormalTok{(}\FunctionTok{paste0}\NormalTok{(path, }\StringTok{"workhistory\_1946.dta"}\NormalTok{)) }\SpecialCharTok{\%\textgreater{}\%}
  \FunctionTok{mutate}\NormalTok{(}\AttributeTok{cohort =} \DecValTok{1946}\NormalTok{)}
\NormalTok{workhistory2 }\OtherTok{\textless{}{-}} \FunctionTok{read\_dta}\NormalTok{(}\FunctionTok{paste0}\NormalTok{(path, }\StringTok{"workhistory\_1958.dta"}\NormalTok{)) }\SpecialCharTok{\%\textgreater{}\%}
  \FunctionTok{mutate}\NormalTok{(}\AttributeTok{cohort =} \DecValTok{1958}\NormalTok{)}
\NormalTok{workhistory3 }\OtherTok{\textless{}{-}} \FunctionTok{read\_dta}\NormalTok{(}\FunctionTok{paste0}\NormalTok{(path, }\StringTok{"workhistory\_1970.dta"}\NormalTok{)) }\SpecialCharTok{\%\textgreater{}\%}
  \FunctionTok{mutate}\NormalTok{(}\AttributeTok{cohort =} \DecValTok{1970}\NormalTok{)}
\end{Highlighting}
\end{Shaded}

The clean social mobility data has about 250,000 rows, while the three
full, raw work history datasets have 5,000-15,000 rows, so I will take
smaller samples of each dataset that I can visually explore in RStudio.

\begin{Shaded}
\begin{Highlighting}[]
\CommentTok{\# Getting samples of each dataset to investigate before looking at summary stats}
\NormalTok{sm\_sample }\OtherTok{\textless{}{-}}\NormalTok{ social\_mobility\_data }\SpecialCharTok{\%\textgreater{}\%} \FunctionTok{slice\_sample}\NormalTok{(}\AttributeTok{n=}\DecValTok{100}\NormalTok{)}
\NormalTok{work\_hist1\_sample }\OtherTok{\textless{}{-}}\NormalTok{ workhistory1 }\SpecialCharTok{\%\textgreater{}\%} \FunctionTok{slice\_sample}\NormalTok{(}\AttributeTok{n=}\DecValTok{100}\NormalTok{)}
\NormalTok{work\_hist2\_sample }\OtherTok{\textless{}{-}}\NormalTok{ workhistory2 }\SpecialCharTok{\%\textgreater{}\%} \FunctionTok{slice\_sample}\NormalTok{(}\AttributeTok{n=}\DecValTok{100}\NormalTok{)}
\NormalTok{work\_hist3\_sample }\OtherTok{\textless{}{-}}\NormalTok{ workhistory3 }\SpecialCharTok{\%\textgreater{}\%} \FunctionTok{slice\_sample}\NormalTok{(}\AttributeTok{n=}\DecValTok{100}\NormalTok{)}
\end{Highlighting}
\end{Shaded}

\subsection{Social Mobility Data}\label{social-mobility-data}

\begin{Shaded}
\begin{Highlighting}[]
\CommentTok{\# Storing the column names and types as a reference}
\NormalTok{social\_mobility\_data }\SpecialCharTok{\%\textgreater{}\%} \FunctionTok{glimpse}\NormalTok{()}
\end{Highlighting}
\end{Shaded}

\begin{verbatim}
## Rows: 248,369
## Columns: 11
## $ .imp     <fct> 0, 0, 0, 0, 0, 0, 0, 0, 0, 0, 0, 0, 0, 0, 0, 0, 0, 0, 0, 0, 0~
## $ .id      <fct> 2, 4, 5, 8, 9, 10, 13, 14, 15, 16, 17, 18, 19, 21, 22, 23, 24~
## $ cohort   <dbl> 1946, 1946, 1946, 1946, 1946, 1946, 1946, 1946, 1946, 1946, 1~
## $ weight   <dbl> 0.4511865, 0.4511865, 0.4511865, 1.8047462, 0.4511865, 0.4511~
## $ id       <dbl> 4600002, 4600004, 4600005, 4600010, 4600011, 4600012, 4600016~
## $ sex      <fct> male, male, male, female, male, male, female, female, female,~
## $ fclass   <fct> Interm, High.sal, Routine, Routine, Interm, Low.sal, Low.sal,~
## $ educ37   <dbl> 7, 7, 1, 3, 4, 3, 5, 4, 4, 4, 6, 2, 1, 3, 1, 1, 1, 2, 1, 7, 4~
## $ class38  <fct> High.sal, NA, Routine, Interm, High.sal, Low.sup, NA, Low.sal~
## $ parttime <dbl> 0, 0, 0, 0, 0, 0, 0, 0, 0, 0, 0, 0, 0, 0, 0, 0, 0, 0, 0, 0, 0~
## $ up       <dbl> 0, NA, NA, NA, NA, NA, NA, NA, NA, NA, NA, NA, NA, NA, NA, NA~
\end{verbatim}

\begin{Shaded}
\begin{Highlighting}[]
\CommentTok{\# Checking the dimensions of the data}
\FunctionTok{dim}\NormalTok{(social\_mobility\_data)}
\end{Highlighting}
\end{Shaded}

\begin{verbatim}
## [1] 248369     11
\end{verbatim}

\subsubsection{Basic data quality
checks}\label{basic-data-quality-checks}

Summary stats on the full data

\begin{Shaded}
\begin{Highlighting}[]
\CommentTok{\# Getting information about the number of missing variables, unique values, basic summary stats}
\FunctionTok{skim}\NormalTok{(social\_mobility\_data) }\SpecialCharTok{\%\textgreater{}\%}
  \FunctionTok{mutate}\NormalTok{(}\FunctionTok{across}\NormalTok{(}\FunctionTok{where}\NormalTok{(is.numeric), }\SpecialCharTok{\textasciitilde{}}\FunctionTok{round}\NormalTok{(., }\DecValTok{2}\NormalTok{)))}
\end{Highlighting}
\end{Shaded}

\begin{longtable}[]{@{}ll@{}}
\caption{Data summary}\tabularnewline
\toprule\noalign{}
\endfirsthead
\endhead
\bottomrule\noalign{}
\endlastfoot
Name & social\_mobility\_data \\
Number of rows & 248369 \\
Number of columns & 11 \\
\_\_\_\_\_\_\_\_\_\_\_\_\_\_\_\_\_\_\_\_\_\_\_ & \\
Column type frequency: & \\
factor & 5 \\
numeric & 6 \\
\_\_\_\_\_\_\_\_\_\_\_\_\_\_\_\_\_\_\_\_\_\_\_\_ & \\
Group variables & None \\
\end{longtable}

\textbf{Variable type: factor}

\begin{longtable}[]{@{}
  >{\raggedright\arraybackslash}p{(\linewidth - 10\tabcolsep) * \real{0.1373}}
  >{\raggedleft\arraybackslash}p{(\linewidth - 10\tabcolsep) * \real{0.0980}}
  >{\raggedleft\arraybackslash}p{(\linewidth - 10\tabcolsep) * \real{0.1373}}
  >{\raggedright\arraybackslash}p{(\linewidth - 10\tabcolsep) * \real{0.0784}}
  >{\raggedleft\arraybackslash}p{(\linewidth - 10\tabcolsep) * \real{0.0882}}
  >{\raggedright\arraybackslash}p{(\linewidth - 10\tabcolsep) * \real{0.4608}}@{}}
\toprule\noalign{}
\begin{minipage}[b]{\linewidth}\raggedright
skim\_variable
\end{minipage} & \begin{minipage}[b]{\linewidth}\raggedleft
n\_missing
\end{minipage} & \begin{minipage}[b]{\linewidth}\raggedleft
complete\_rate
\end{minipage} & \begin{minipage}[b]{\linewidth}\raggedright
ordered
\end{minipage} & \begin{minipage}[b]{\linewidth}\raggedleft
n\_unique
\end{minipage} & \begin{minipage}[b]{\linewidth}\raggedright
top\_counts
\end{minipage} \\
\midrule\noalign{}
\endhead
\bottomrule\noalign{}
\endlastfoot
.imp & 0 & 1.00 & FALSE & 11 & 0: 22579, 1: 22579, 2: 22579, 3: 22579 \\
.id & 0 & 1.00 & FALSE & 22579 & 10: 11, 100: 11, 100: 11, 100: 11 \\
sex & 0 & 1.00 & FALSE & 2 & mal: 171512, fem: 76857 \\
fclass & 2812 & 0.99 & FALSE & 7 & Rou: 62556, Low: 40657, Low: 38920,
Sem: 31499 \\
class38 & 7099 & 0.97 & FALSE & 7 & Low: 55670, Int: 37531, Hig: 36356,
Sem: 33915 \\
\end{longtable}

\textbf{Variable type: numeric}

\begin{longtable}[]{@{}
  >{\raggedright\arraybackslash}p{(\linewidth - 20\tabcolsep) * \real{0.1273}}
  >{\raggedleft\arraybackslash}p{(\linewidth - 20\tabcolsep) * \real{0.0909}}
  >{\raggedleft\arraybackslash}p{(\linewidth - 20\tabcolsep) * \real{0.1273}}
  >{\raggedleft\arraybackslash}p{(\linewidth - 20\tabcolsep) * \real{0.1000}}
  >{\raggedleft\arraybackslash}p{(\linewidth - 20\tabcolsep) * \real{0.0909}}
  >{\raggedleft\arraybackslash}p{(\linewidth - 20\tabcolsep) * \real{0.1000}}
  >{\raggedleft\arraybackslash}p{(\linewidth - 20\tabcolsep) * \real{0.0727}}
  >{\raggedleft\arraybackslash}p{(\linewidth - 20\tabcolsep) * \real{0.0727}}
  >{\raggedleft\arraybackslash}p{(\linewidth - 20\tabcolsep) * \real{0.0727}}
  >{\raggedleft\arraybackslash}p{(\linewidth - 20\tabcolsep) * \real{0.0909}}
  >{\raggedright\arraybackslash}p{(\linewidth - 20\tabcolsep) * \real{0.0545}}@{}}
\toprule\noalign{}
\begin{minipage}[b]{\linewidth}\raggedright
skim\_variable
\end{minipage} & \begin{minipage}[b]{\linewidth}\raggedleft
n\_missing
\end{minipage} & \begin{minipage}[b]{\linewidth}\raggedleft
complete\_rate
\end{minipage} & \begin{minipage}[b]{\linewidth}\raggedleft
mean
\end{minipage} & \begin{minipage}[b]{\linewidth}\raggedleft
sd
\end{minipage} & \begin{minipage}[b]{\linewidth}\raggedleft
p0
\end{minipage} & \begin{minipage}[b]{\linewidth}\raggedleft
p25
\end{minipage} & \begin{minipage}[b]{\linewidth}\raggedleft
p50
\end{minipage} & \begin{minipage}[b]{\linewidth}\raggedleft
p75
\end{minipage} & \begin{minipage}[b]{\linewidth}\raggedleft
p100
\end{minipage} & \begin{minipage}[b]{\linewidth}\raggedright
hist
\end{minipage} \\
\midrule\noalign{}
\endhead
\bottomrule\noalign{}
\endlastfoot
cohort & 0 & 1.00 & 1960.66 & 8.27 & 1946.00 & 1958 & 1958 & 1970 &
1970.0 & ▂▁▇▁▆ \\
weight & 0 & 1.00 & 1.00 & 0.26 & 0.45 & 1 & 1 & 1 & 1.8 & ▁▁▇▁▁ \\
id & 0 & 1.00 & 6071657.16 & 827557.30 & 4600002.00 & 5802954 & 5810508
& 7003984 & 7012132.0 & ▂▁▇▁▆ \\
educ37 & 84 & 1.00 & 3.78 & 2.12 & 1.00 & 2 & 4 & 6 & 8.0 & ▇▆▆▃▅ \\
parttime & 1859 & 0.99 & 0.05 & 0.23 & 0.00 & 0 & 0 & 0 & 1.0 & ▇▁▁▁▁ \\
up & 248368 & 0.00 & 0.00 & NA & 0.00 & 0 & 0 & 0 & 0.0 & ▁▁▇▁▁ \\
\end{longtable}

\subsubsection{Data cleaning}\label{data-cleaning}

\paragraph{Restricting to only observations with full
data}\label{restricting-to-only-observations-with-full-data}

The \texttt{.imp} variable is equal to \(0\) when the row of data was
original and was not multiply imputed and is not equal to \(0\) when
some of the data was multiply imputed to deal with missing data. There
are many strange properties of the full data, like two ID variables that
frequently appear multiple times in the data. By restricting the sample
to only those that have full data, it is easier to see the properties of
the data and to conduct baseline analyses. The selection bias that such
a filtering of the data might induce can be investigated later.

\paragraph{Recoding data types}\label{recoding-data-types}

It is unclear as of now what the difference between \texttt{id} and
\texttt{.id} are. However, \texttt{id} and \texttt{cohort} should
clearly be unordered, categorical variables.

\begin{Shaded}
\begin{Highlighting}[]
\NormalTok{sm\_data1 }\OtherTok{\textless{}{-}}\NormalTok{ social\_mobility\_data }\SpecialCharTok{\%\textgreater{}\%}
  \FunctionTok{filter}\NormalTok{(.imp }\SpecialCharTok{==} \DecValTok{0}\NormalTok{) }\SpecialCharTok{\%\textgreater{}\%}
  \FunctionTok{mutate}\NormalTok{(}\AttributeTok{id =} \FunctionTok{as.factor}\NormalTok{(id),}
         \AttributeTok{cohort =} \FunctionTok{as.factor}\NormalTok{(cohort))}
\end{Highlighting}
\end{Shaded}

\begin{Shaded}
\begin{Highlighting}[]
\CommentTok{\# Checking the dimensions of the data}
\FunctionTok{dim}\NormalTok{(sm\_data1)}
\end{Highlighting}
\end{Shaded}

\begin{verbatim}
## [1] 22579    11
\end{verbatim}

\begin{Shaded}
\begin{Highlighting}[]
\CommentTok{\# Getting information about the number of missing variables, unique values, basic summary stats}
\FunctionTok{skim}\NormalTok{(sm\_data1) }\SpecialCharTok{\%\textgreater{}\%}
  \FunctionTok{mutate}\NormalTok{(}\FunctionTok{across}\NormalTok{(}\FunctionTok{where}\NormalTok{(is.numeric), }\SpecialCharTok{\textasciitilde{}}\FunctionTok{round}\NormalTok{(., }\DecValTok{2}\NormalTok{)))}
\end{Highlighting}
\end{Shaded}

\begin{longtable}[]{@{}ll@{}}
\caption{Data summary}\tabularnewline
\toprule\noalign{}
\endfirsthead
\endhead
\bottomrule\noalign{}
\endlastfoot
Name & sm\_data1 \\
Number of rows & 22579 \\
Number of columns & 11 \\
\_\_\_\_\_\_\_\_\_\_\_\_\_\_\_\_\_\_\_\_\_\_\_ & \\
Column type frequency: & \\
factor & 7 \\
numeric & 4 \\
\_\_\_\_\_\_\_\_\_\_\_\_\_\_\_\_\_\_\_\_\_\_\_\_ & \\
Group variables & None \\
\end{longtable}

\textbf{Variable type: factor}

\begin{longtable}[]{@{}
  >{\raggedright\arraybackslash}p{(\linewidth - 10\tabcolsep) * \real{0.1429}}
  >{\raggedleft\arraybackslash}p{(\linewidth - 10\tabcolsep) * \real{0.1020}}
  >{\raggedleft\arraybackslash}p{(\linewidth - 10\tabcolsep) * \real{0.1429}}
  >{\raggedright\arraybackslash}p{(\linewidth - 10\tabcolsep) * \real{0.0816}}
  >{\raggedleft\arraybackslash}p{(\linewidth - 10\tabcolsep) * \real{0.0918}}
  >{\raggedright\arraybackslash}p{(\linewidth - 10\tabcolsep) * \real{0.4388}}@{}}
\toprule\noalign{}
\begin{minipage}[b]{\linewidth}\raggedright
skim\_variable
\end{minipage} & \begin{minipage}[b]{\linewidth}\raggedleft
n\_missing
\end{minipage} & \begin{minipage}[b]{\linewidth}\raggedleft
complete\_rate
\end{minipage} & \begin{minipage}[b]{\linewidth}\raggedright
ordered
\end{minipage} & \begin{minipage}[b]{\linewidth}\raggedleft
n\_unique
\end{minipage} & \begin{minipage}[b]{\linewidth}\raggedright
top\_counts
\end{minipage} \\
\midrule\noalign{}
\endhead
\bottomrule\noalign{}
\endlastfoot
.imp & 0 & 1.00 & FALSE & 1 & 0: 22579, 1: 0, 2: 0, 3: 0 \\
.id & 0 & 1.00 & FALSE & 22579 & 10: 1, 100: 1, 100: 1, 100: 1 \\
cohort & 0 & 1.00 & FALSE & 3 & 195: 10754, 197: 8411, 194: 3414 \\
id & 0 & 1.00 & FALSE & 22579 & 460: 1, 460: 1, 460: 1, 460: 1 \\
sex & 0 & 1.00 & FALSE & 2 & mal: 15592, fem: 6987 \\
fclass & 2812 & 0.88 & FALSE & 7 & Rou: 5010, Low: 3298, Low: 3145, Sem:
2507 \\
class38 & 7099 & 0.69 & FALSE & 7 & Low: 3723, Hig: 2523, Int: 2442,
Sem: 1979 \\
\end{longtable}

\textbf{Variable type: numeric}

\begin{longtable}[]{@{}
  >{\raggedright\arraybackslash}p{(\linewidth - 20\tabcolsep) * \real{0.1842}}
  >{\raggedleft\arraybackslash}p{(\linewidth - 20\tabcolsep) * \real{0.1316}}
  >{\raggedleft\arraybackslash}p{(\linewidth - 20\tabcolsep) * \real{0.1842}}
  >{\raggedleft\arraybackslash}p{(\linewidth - 20\tabcolsep) * \real{0.0658}}
  >{\raggedleft\arraybackslash}p{(\linewidth - 20\tabcolsep) * \real{0.0658}}
  >{\raggedleft\arraybackslash}p{(\linewidth - 20\tabcolsep) * \real{0.0658}}
  >{\raggedleft\arraybackslash}p{(\linewidth - 20\tabcolsep) * \real{0.0526}}
  >{\raggedleft\arraybackslash}p{(\linewidth - 20\tabcolsep) * \real{0.0526}}
  >{\raggedleft\arraybackslash}p{(\linewidth - 20\tabcolsep) * \real{0.0526}}
  >{\raggedleft\arraybackslash}p{(\linewidth - 20\tabcolsep) * \real{0.0658}}
  >{\raggedright\arraybackslash}p{(\linewidth - 20\tabcolsep) * \real{0.0789}}@{}}
\toprule\noalign{}
\begin{minipage}[b]{\linewidth}\raggedright
skim\_variable
\end{minipage} & \begin{minipage}[b]{\linewidth}\raggedleft
n\_missing
\end{minipage} & \begin{minipage}[b]{\linewidth}\raggedleft
complete\_rate
\end{minipage} & \begin{minipage}[b]{\linewidth}\raggedleft
mean
\end{minipage} & \begin{minipage}[b]{\linewidth}\raggedleft
sd
\end{minipage} & \begin{minipage}[b]{\linewidth}\raggedleft
p0
\end{minipage} & \begin{minipage}[b]{\linewidth}\raggedleft
p25
\end{minipage} & \begin{minipage}[b]{\linewidth}\raggedleft
p50
\end{minipage} & \begin{minipage}[b]{\linewidth}\raggedleft
p75
\end{minipage} & \begin{minipage}[b]{\linewidth}\raggedleft
p100
\end{minipage} & \begin{minipage}[b]{\linewidth}\raggedright
hist
\end{minipage} \\
\midrule\noalign{}
\endhead
\bottomrule\noalign{}
\endlastfoot
weight & 0 & 1.00 & 1.00 & 0.26 & 0.45 & 1 & 1 & 1 & 1.8 & ▁▁▇▁▁ \\
educ37 & 84 & 1.00 & 3.78 & 2.12 & 1.00 & 2 & 4 & 6 & 8.0 & ▇▆▆▃▅ \\
parttime & 169 & 0.99 & 0.05 & 0.23 & 0.00 & 0 & 0 & 0 & 1.0 & ▇▁▁▁▁ \\
up & 22578 & 0.00 & 0.00 & NA & 0.00 & 0 & 0 & 0 & 0.0 & ▁▁▇▁▁ \\
\end{longtable}

It is still unclear what the relationship between the two ID variables
is, and which one should be used to map to the IQ testing data. It is
also unclear what the \texttt{weight}, \texttt{parttime}, and
\texttt{up} variables are meant to map to.

\subsubsection{Variable definitions}\label{variable-definitions}

As of this point, these are my working definitions of each of the
variables:

\begin{itemize}
\tightlist
\item
  \texttt{id}, \texttt{.id}: The unit of observation/person in the data.
  Unclear as of now which variable is the relevant one.
\item
  \texttt{educ37}: Ordinal scale from 1-8 indicating the unit's level of
  education at age 37. The mediating variable in Kuha et al (2021), and
  one of our variables of interest here.
\item
  \texttt{fclass}: The unit's father's social class, as described in
  Kuha et al (2021). The ``treatment'' variable (though it does not
  exactly admit a causal interpretation here).
\item
  \texttt{class38}: The unit's social class at age 38. The outcome
  variable.
\item
  \texttt{.imp}: Indicator of whether the data was imputed or not.
\item
  \texttt{cohort}: The year of birth for the unit being measured. There
  are three cohorts, each 12 years apart, in the data.
\item
  \texttt{sex}: The sex of the unit.
\item
  \texttt{weight}: This is a survey weight for units from the 1946
  cohort, described in Section 2 of Kuha et al (2021).
\item
  \texttt{parttime}, \texttt{up}: Unknown for now. \texttt{weight} might
  be some kind of adjustment for sample over/under representation of the
  true population.
\end{itemize}

It is also helpful to describe what each of the 7 levels that the social
class variables can take are. This will also help justify the
simplifying assumption I will make below for visualization purposes to
interpret the class levels at ordered, which is not strictly accurate.
From Kuha et al (2021):

\begin{verbatim}
1.  Higher managers and professionals

2.  Lower managers and professionals

3.  Intermediate occupations

4.  Small employers and own account workers

5.  Lower supervisory and technical occupations

6.  Semiroutine occupations

7.  Routine occupations
\end{verbatim}

\subsubsection{Relationships between
variables}\label{relationships-between-variables}

\paragraph{Evolution of the treatment, mediator, and outcome variable
distributions over
time}\label{evolution-of-the-treatment-mediator-and-outcome-variable-distributions-over-time}

\begin{Shaded}
\begin{Highlighting}[]
\CommentTok{\# I pretend that the "roughly ordinal" ordering of the social class variables}
\CommentTok{\# are strict orderings for visualization purposes}

\NormalTok{sm\_data2 }\OtherTok{\textless{}{-}}\NormalTok{ sm\_data1 }\SpecialCharTok{\%\textgreater{}\%}
  \FunctionTok{mutate}\NormalTok{(}\AttributeTok{fclass =} \FunctionTok{as.numeric}\NormalTok{(fclass), }
         \AttributeTok{class38 =} \FunctionTok{as.numeric}\NormalTok{(class38))}

\NormalTok{sm\_data2 }\SpecialCharTok{\%\textgreater{}\%}
  \FunctionTok{pivot\_longer}\NormalTok{(}\AttributeTok{cols =} \FunctionTok{c}\NormalTok{(educ37, class38, fclass), }\AttributeTok{names\_to =} \StringTok{"variable"}\NormalTok{, }\AttributeTok{values\_to =} \StringTok{"value"}\NormalTok{) }\SpecialCharTok{\%\textgreater{}\%}
  \FunctionTok{ggplot}\NormalTok{(}\FunctionTok{aes}\NormalTok{(}\AttributeTok{x =}\NormalTok{ value, }\AttributeTok{fill =}\NormalTok{ cohort)) }\SpecialCharTok{+}
  \FunctionTok{geom\_boxplot}\NormalTok{(}\AttributeTok{alpha =} \FloatTok{0.5}\NormalTok{) }\SpecialCharTok{+} 
  \FunctionTok{facet\_wrap}\NormalTok{(}\SpecialCharTok{\textasciitilde{}}\NormalTok{variable, }\AttributeTok{scales =} \StringTok{"free"}\NormalTok{) }\SpecialCharTok{+} 
  \FunctionTok{guides}\NormalTok{(}\AttributeTok{fill =} \FunctionTok{guide\_legend}\NormalTok{(}\AttributeTok{reverse =} \ConstantTok{TRUE}\NormalTok{)) }\SpecialCharTok{+} 
  \FunctionTok{theme}\NormalTok{(}\AttributeTok{axis.text.x =} \FunctionTok{element\_blank}\NormalTok{(),}
        \AttributeTok{axis.title.x =} \FunctionTok{element\_blank}\NormalTok{())}
\end{Highlighting}
\end{Shaded}

\pandocbounded{\includegraphics[keepaspectratio]{social_mobility_data_eda_files/figure-latex/unnamed-chunk-9-1.pdf}}

As expected, level of education seems to increase over time, while
social class and father's social class seem to remain relatively stable
(with a slight change/ approximate ``increase'' in father's social
class).

\paragraph{Associations of treatment, mediator,
outcome}\label{associations-of-treatment-mediator-outcome}

Running a series of basic regressions on the association between all
three pairs of variables. Once again, social class is treated as roughly
ordinal, where lower values are associated with social classes that I
would expect to be roughly seen as ``more advantaged''.

\begin{Shaded}
\begin{Highlighting}[]
\CommentTok{\# Creating a function to fit a linear regression and visualize + summarize the relationship}
\NormalTok{fit\_lm }\OtherTok{\textless{}{-}} \ControlFlowTok{function}\NormalTok{(x, y) \{}
\NormalTok{  model }\OtherTok{\textless{}{-}} \FunctionTok{lm}\NormalTok{(}\FunctionTok{reformulate}\NormalTok{(x, y), }\AttributeTok{data =}\NormalTok{ sm\_data2)}
  
  \FunctionTok{print}\NormalTok{(}\FunctionTok{summary}\NormalTok{(model))}
  
  \FunctionTok{ggplot}\NormalTok{(sm\_data2, }\FunctionTok{aes}\NormalTok{(}\AttributeTok{x =}\NormalTok{ .data[[x]], }\AttributeTok{y =}\NormalTok{ .data[[y]])) }\SpecialCharTok{+}
    \FunctionTok{geom\_smooth}\NormalTok{(}\AttributeTok{method =} \StringTok{"lm"}\NormalTok{, }\AttributeTok{se =} \ConstantTok{TRUE}\NormalTok{) }\SpecialCharTok{+}
    \FunctionTok{scale\_x\_continuous}\NormalTok{(}\AttributeTok{limits =} \FunctionTok{c}\NormalTok{(}\FunctionTok{min}\NormalTok{(sm\_data2[[x]], }\AttributeTok{na.rm =} \ConstantTok{TRUE}\NormalTok{), }\FunctionTok{max}\NormalTok{(sm\_data2[[x]], }\AttributeTok{na.rm =} \ConstantTok{TRUE}\NormalTok{))) }\SpecialCharTok{+}
    \FunctionTok{scale\_y\_continuous}\NormalTok{(}\AttributeTok{limits =} \FunctionTok{c}\NormalTok{(}\FunctionTok{min}\NormalTok{(sm\_data2[[y]], }\AttributeTok{na.rm =} \ConstantTok{TRUE}\NormalTok{), }\FunctionTok{max}\NormalTok{(sm\_data2[[y]], }\AttributeTok{na.rm =} \ConstantTok{TRUE}\NormalTok{))) }\SpecialCharTok{+}
    \FunctionTok{labs}\NormalTok{(}
      \AttributeTok{title =} \StringTok{""}\NormalTok{,}
      \AttributeTok{x =}\NormalTok{ x,}
      \AttributeTok{y =}\NormalTok{ y}
\NormalTok{    ) }\SpecialCharTok{+}
    \FunctionTok{theme\_classic}\NormalTok{()}
\NormalTok{\}}
\end{Highlighting}
\end{Shaded}

\begin{Shaded}
\begin{Highlighting}[]
\CommentTok{\# Father\textquotesingle{}s social class {-}\textgreater{} education}
\FunctionTok{fit\_lm}\NormalTok{(}\StringTok{"fclass"}\NormalTok{, }\StringTok{"educ37"}\NormalTok{)}
\end{Highlighting}
\end{Shaded}

\begin{verbatim}
## 
## Call:
## lm(formula = reformulate(x, y), data = sm_data2)
## 
## Residuals:
##     Min      1Q  Median      3Q     Max 
## -4.0342 -1.6867 -0.0342  1.6608  5.0508 
## 
## Coefficients:
##              Estimate Std. Error t value Pr(>|t|)    
## (Intercept)  5.381684   0.033659  159.89   <2e-16 ***
## fclass      -0.347499   0.006841  -50.79   <2e-16 ***
## ---
## Signif. codes:  0 '***' 0.001 '**' 0.01 '*' 0.05 '.' 0.1 ' ' 1
## 
## Residual standard error: 1.989 on 19721 degrees of freedom
##   (2856 observations deleted due to missingness)
## Multiple R-squared:  0.1157, Adjusted R-squared:  0.1156 
## F-statistic:  2580 on 1 and 19721 DF,  p-value: < 2.2e-16
\end{verbatim}

\begin{verbatim}
## `geom_smooth()` using formula = 'y ~ x'
\end{verbatim}

\pandocbounded{\includegraphics[keepaspectratio]{social_mobility_data_eda_files/figure-latex/unnamed-chunk-11-1.pdf}}

\begin{Shaded}
\begin{Highlighting}[]
\CommentTok{\# Father\textquotesingle{}s social class {-}\textgreater{} unit social class}
\FunctionTok{fit\_lm}\NormalTok{(}\StringTok{"fclass"}\NormalTok{, }\StringTok{"class38"}\NormalTok{)}
\end{Highlighting}
\end{Shaded}

\begin{verbatim}
## 
## Call:
## lm(formula = reformulate(x, y), data = sm_data2)
## 
## Residuals:
##     Min      1Q  Median      3Q     Max 
## -3.3369 -1.5436 -0.4402  1.6631  4.4564 
## 
## Coefficients:
##             Estimate Std. Error t value Pr(>|t|)    
## (Intercept) 2.244664   0.038085   58.94   <2e-16 ***
## fclass      0.298891   0.007875   37.96   <2e-16 ***
## ---
## Signif. codes:  0 '***' 0.001 '**' 0.01 '*' 0.05 '.' 0.1 ' ' 1
## 
## Residual standard error: 1.92 on 13797 degrees of freedom
##   (8780 observations deleted due to missingness)
## Multiple R-squared:  0.09455,    Adjusted R-squared:  0.09448 
## F-statistic:  1441 on 1 and 13797 DF,  p-value: < 2.2e-16
\end{verbatim}

\begin{verbatim}
## `geom_smooth()` using formula = 'y ~ x'
\end{verbatim}

\pandocbounded{\includegraphics[keepaspectratio]{social_mobility_data_eda_files/figure-latex/unnamed-chunk-12-1.pdf}}

\begin{Shaded}
\begin{Highlighting}[]
\CommentTok{\# Education {-}\textgreater{} class38}
\FunctionTok{fit\_lm}\NormalTok{(}\StringTok{"educ37"}\NormalTok{, }\StringTok{"class38"}\NormalTok{)}
\end{Highlighting}
\end{Shaded}

\begin{verbatim}
## 
## Call:
## lm(formula = reformulate(x, y), data = sm_data2)
## 
## Residuals:
##    Min     1Q Median     3Q    Max 
## -3.986 -1.148 -0.148  1.433  5.325 
## 
## Coefficients:
##              Estimate Std. Error t value Pr(>|t|)    
## (Intercept)  5.459063   0.030269  180.35   <2e-16 ***
## educ37      -0.473011   0.006697  -70.63   <2e-16 ***
## ---
## Signif. codes:  0 '***' 0.001 '**' 0.01 '*' 0.05 '.' 0.1 ' ' 1
## 
## Residual standard error: 1.758 on 15446 degrees of freedom
##   (7131 observations deleted due to missingness)
## Multiple R-squared:  0.2441, Adjusted R-squared:  0.2441 
## F-statistic:  4989 on 1 and 15446 DF,  p-value: < 2.2e-16
\end{verbatim}

\begin{verbatim}
## `geom_smooth()` using formula = 'y ~ x'
\end{verbatim}

\pandocbounded{\includegraphics[keepaspectratio]{social_mobility_data_eda_files/figure-latex/unnamed-chunk-13-1.pdf}}

\subsection{Work history data}\label{work-history-data}

\begin{Shaded}
\begin{Highlighting}[]
\CommentTok{\# Checking the dimensions of the data}
\FunctionTok{dim}\NormalTok{(workhistory1)}
\end{Highlighting}
\end{Shaded}

\begin{verbatim}
## [1] 4506  144
\end{verbatim}

\begin{Shaded}
\begin{Highlighting}[]
\FunctionTok{dim}\NormalTok{(workhistory2)}
\end{Highlighting}
\end{Shaded}

\begin{verbatim}
## [1] 14313   232
\end{verbatim}

\begin{Shaded}
\begin{Highlighting}[]
\FunctionTok{dim}\NormalTok{(workhistory3)}
\end{Highlighting}
\end{Shaded}

\begin{verbatim}
## [1] 12132   170
\end{verbatim}

A cursory look at the data and the number of variables shows that these
data contain detailed information about the evolution of relevant
unit-level characteristics: social class, education, and occupation. Of
these variables, the one that is of immediate interest is the
\texttt{IQ} variable as well as the other variables in the dataset that
are relevant to the original analysis in Kuha et al (2021) and are
included in the cleaned dataset.

\begin{Shaded}
\begin{Highlighting}[]
\CommentTok{\# Creating a combined dataset}
\NormalTok{wk\_data1 }\OtherTok{\textless{}{-}} \FunctionTok{bind\_rows}\NormalTok{(workhistory1 }\SpecialCharTok{\%\textgreater{}\%} \FunctionTok{mutate}\NormalTok{(}\AttributeTok{cohort =} \StringTok{"1946"}\NormalTok{,}
                                              \AttributeTok{ed37 =}\NormalTok{ educ), }
\NormalTok{                      workhistory2 }\SpecialCharTok{\%\textgreater{}\%} \FunctionTok{mutate}\NormalTok{(}\AttributeTok{cohort =} \StringTok{"1958"}\NormalTok{),}
\NormalTok{                      workhistory3 }\SpecialCharTok{\%\textgreater{}\%} \FunctionTok{mutate}\NormalTok{(}\AttributeTok{cohort =} \StringTok{"1970"}\NormalTok{)) }\SpecialCharTok{\%\textgreater{}\%}
  \FunctionTok{select}\NormalTok{(id, sex, fclass, ed37, class38, IQ, cohort) }\SpecialCharTok{\%\textgreater{}\%}
  \FunctionTok{mutate}\NormalTok{(}\AttributeTok{fclass =} \FunctionTok{as.factor}\NormalTok{(fclass),}
         \AttributeTok{ed37 =} \FunctionTok{as.numeric}\NormalTok{(ed37),}
         \AttributeTok{class38 =} \FunctionTok{as.factor}\NormalTok{(class38))}
\end{Highlighting}
\end{Shaded}

\begin{Shaded}
\begin{Highlighting}[]
\CommentTok{\# Basic summary statistics}
\FunctionTok{skim}\NormalTok{(wk\_data1) }\SpecialCharTok{\%\textgreater{}\%}
  \FunctionTok{mutate}\NormalTok{(}\FunctionTok{across}\NormalTok{(}\FunctionTok{where}\NormalTok{(is.numeric), }\SpecialCharTok{\textasciitilde{}}\FunctionTok{round}\NormalTok{(., }\DecValTok{2}\NormalTok{)))}
\end{Highlighting}
\end{Shaded}

\begin{longtable}[]{@{}ll@{}}
\caption{Data summary}\tabularnewline
\toprule\noalign{}
\endfirsthead
\endhead
\bottomrule\noalign{}
\endlastfoot
Name & wk\_data1 \\
Number of rows & 30951 \\
Number of columns & 7 \\
\_\_\_\_\_\_\_\_\_\_\_\_\_\_\_\_\_\_\_\_\_\_\_ & \\
Column type frequency: & \\
character & 1 \\
factor & 2 \\
numeric & 4 \\
\_\_\_\_\_\_\_\_\_\_\_\_\_\_\_\_\_\_\_\_\_\_\_\_ & \\
Group variables & None \\
\end{longtable}

\textbf{Variable type: character}

\begin{longtable}[]{@{}
  >{\raggedright\arraybackslash}p{(\linewidth - 14\tabcolsep) * \real{0.1944}}
  >{\raggedleft\arraybackslash}p{(\linewidth - 14\tabcolsep) * \real{0.1389}}
  >{\raggedleft\arraybackslash}p{(\linewidth - 14\tabcolsep) * \real{0.1944}}
  >{\raggedleft\arraybackslash}p{(\linewidth - 14\tabcolsep) * \real{0.0556}}
  >{\raggedleft\arraybackslash}p{(\linewidth - 14\tabcolsep) * \real{0.0556}}
  >{\raggedleft\arraybackslash}p{(\linewidth - 14\tabcolsep) * \real{0.0833}}
  >{\raggedleft\arraybackslash}p{(\linewidth - 14\tabcolsep) * \real{0.1250}}
  >{\raggedleft\arraybackslash}p{(\linewidth - 14\tabcolsep) * \real{0.1528}}@{}}
\toprule\noalign{}
\begin{minipage}[b]{\linewidth}\raggedright
skim\_variable
\end{minipage} & \begin{minipage}[b]{\linewidth}\raggedleft
n\_missing
\end{minipage} & \begin{minipage}[b]{\linewidth}\raggedleft
complete\_rate
\end{minipage} & \begin{minipage}[b]{\linewidth}\raggedleft
min
\end{minipage} & \begin{minipage}[b]{\linewidth}\raggedleft
max
\end{minipage} & \begin{minipage}[b]{\linewidth}\raggedleft
empty
\end{minipage} & \begin{minipage}[b]{\linewidth}\raggedleft
n\_unique
\end{minipage} & \begin{minipage}[b]{\linewidth}\raggedleft
whitespace
\end{minipage} \\
\midrule\noalign{}
\endhead
\bottomrule\noalign{}
\endlastfoot
cohort & 0 & 1 & 4 & 4 & 0 & 3 & 0 \\
\end{longtable}

\textbf{Variable type: factor}

\begin{longtable}[]{@{}
  >{\raggedright\arraybackslash}p{(\linewidth - 10\tabcolsep) * \real{0.1538}}
  >{\raggedleft\arraybackslash}p{(\linewidth - 10\tabcolsep) * \real{0.1099}}
  >{\raggedleft\arraybackslash}p{(\linewidth - 10\tabcolsep) * \real{0.1538}}
  >{\raggedright\arraybackslash}p{(\linewidth - 10\tabcolsep) * \real{0.0879}}
  >{\raggedleft\arraybackslash}p{(\linewidth - 10\tabcolsep) * \real{0.0989}}
  >{\raggedright\arraybackslash}p{(\linewidth - 10\tabcolsep) * \real{0.3956}}@{}}
\toprule\noalign{}
\begin{minipage}[b]{\linewidth}\raggedright
skim\_variable
\end{minipage} & \begin{minipage}[b]{\linewidth}\raggedleft
n\_missing
\end{minipage} & \begin{minipage}[b]{\linewidth}\raggedleft
complete\_rate
\end{minipage} & \begin{minipage}[b]{\linewidth}\raggedright
ordered
\end{minipage} & \begin{minipage}[b]{\linewidth}\raggedleft
n\_unique
\end{minipage} & \begin{minipage}[b]{\linewidth}\raggedright
top\_counts
\end{minipage} \\
\midrule\noalign{}
\endhead
\bottomrule\noalign{}
\endlastfoot
fclass & 3788 & 0.88 & FALSE & 7 & 7: 6852, 2: 4525, 5: 4395, 6: 3463 \\
class38 & 0 & 1.00 & FALSE & 10 & -1: 7679, 2: 5290, 3: 4154, 6: 3715 \\
\end{longtable}

\textbf{Variable type: numeric}

\begin{longtable}[]{@{}
  >{\raggedright\arraybackslash}p{(\linewidth - 20\tabcolsep) * \real{0.1414}}
  >{\raggedleft\arraybackslash}p{(\linewidth - 20\tabcolsep) * \real{0.1010}}
  >{\raggedleft\arraybackslash}p{(\linewidth - 20\tabcolsep) * \real{0.1414}}
  >{\raggedleft\arraybackslash}p{(\linewidth - 20\tabcolsep) * \real{0.0808}}
  >{\raggedleft\arraybackslash}p{(\linewidth - 20\tabcolsep) * \real{0.0808}}
  >{\raggedleft\arraybackslash}p{(\linewidth - 20\tabcolsep) * \real{0.0606}}
  >{\raggedleft\arraybackslash}p{(\linewidth - 20\tabcolsep) * \real{0.0808}}
  >{\raggedleft\arraybackslash}p{(\linewidth - 20\tabcolsep) * \real{0.0808}}
  >{\raggedleft\arraybackslash}p{(\linewidth - 20\tabcolsep) * \real{0.0808}}
  >{\raggedleft\arraybackslash}p{(\linewidth - 20\tabcolsep) * \real{0.0909}}
  >{\raggedright\arraybackslash}p{(\linewidth - 20\tabcolsep) * \real{0.0606}}@{}}
\toprule\noalign{}
\begin{minipage}[b]{\linewidth}\raggedright
skim\_variable
\end{minipage} & \begin{minipage}[b]{\linewidth}\raggedleft
n\_missing
\end{minipage} & \begin{minipage}[b]{\linewidth}\raggedleft
complete\_rate
\end{minipage} & \begin{minipage}[b]{\linewidth}\raggedleft
mean
\end{minipage} & \begin{minipage}[b]{\linewidth}\raggedleft
sd
\end{minipage} & \begin{minipage}[b]{\linewidth}\raggedleft
p0
\end{minipage} & \begin{minipage}[b]{\linewidth}\raggedleft
p25
\end{minipage} & \begin{minipage}[b]{\linewidth}\raggedleft
p50
\end{minipage} & \begin{minipage}[b]{\linewidth}\raggedleft
p75
\end{minipage} & \begin{minipage}[b]{\linewidth}\raggedleft
p100
\end{minipage} & \begin{minipage}[b]{\linewidth}\raggedright
hist
\end{minipage} \\
\midrule\noalign{}
\endhead
\bottomrule\noalign{}
\endlastfoot
id & 12132 & 0.61 & 6100.95 & 4151.55 & 1.00 & 2578.00 & 5124.00 &
9622.50 & 14346.00 & ▇▇▅▅▅ \\
sex & 0 & 1.00 & 0.50 & 0.50 & 0.00 & 0.00 & 0.00 & 1.00 & 1.00 &
▇▁▁▁▇ \\
ed37 & 112 & 1.00 & 3.76 & 2.09 & 1.00 & 2.00 & 3.00 & 6.00 & 8.00 &
▇▆▆▃▃ \\
IQ & 10186 & 0.67 & 0.03 & 0.99 & -3.66 & -0.67 & 0.08 & 0.77 & 3.37 &
▁▃▇▅▁ \\
\end{longtable}

\subsubsection{Comparison of the three
datasets}\label{comparison-of-the-three-datasets}

\begin{Shaded}
\begin{Highlighting}[]
\CommentTok{\# Checking overlap in variables between the three cohort datasets}
\FunctionTok{intersect}\NormalTok{(}\FunctionTok{intersect}\NormalTok{(}\FunctionTok{colnames}\NormalTok{(workhistory1), }\FunctionTok{colnames}\NormalTok{(workhistory2)),  }
          \FunctionTok{colnames}\NormalTok{(workhistory3))}
\end{Highlighting}
\end{Shaded}

\begin{verbatim}
##  [1] "sex"     "part16"  "part17"  "part18"  "part19"  "part20"  "part21" 
##  [8] "part22"  "part23"  "part24"  "part25"  "part26"  "part27"  "part28" 
## [15] "part29"  "part30"  "part31"  "part32"  "part33"  "part34"  "part35" 
## [22] "part36"  "part37"  "part38"  "nssec16" "nssec17" "nssec18" "nssec19"
## [29] "nssec20" "nssec21" "nssec22" "nssec23" "nssec24" "nssec25" "nssec26"
## [36] "nssec27" "nssec28" "nssec29" "nssec30" "nssec31" "nssec32" "nssec33"
## [43] "nssec34" "nssec35" "nssec36" "nssec37" "nssec38" "class16" "class17"
## [50] "class18" "class19" "class20" "class21" "class22" "class23" "class24"
## [57] "class25" "class26" "class27" "class28" "class29" "class30" "class31"
## [64] "class32" "class33" "class34" "class35" "class36" "class37" "class38"
## [71] "occ16"   "occ17"   "occ18"   "occ19"   "occ20"   "occ21"   "occ22"  
## [78] "occ23"   "occ24"   "occ25"   "occ26"   "occ27"   "occ28"   "occ29"  
## [85] "occ30"   "occ31"   "occ32"   "occ33"   "occ34"   "occ35"   "occ36"  
## [92] "occ37"   "occ38"   "fclass"  "IQ"      "ed426"   "cohort"
\end{verbatim}

An immediate concern is that the 1970 data does not seem to have the
same \texttt{id} variable as the 1946 and 1958 cohort data. Instead,
there is a \texttt{bcsid} variable that does not have a clear analogue
in the cleaned data or the other two data files. However, depending on
well the first two work history datasets match with the cleaned social
mobility data, it is possible that only the work history datasets could
be used in the analysis. Instead of \texttt{ed37}, the 1946 cohort data
has a \texttt{educ} variable.

\subsubsection{Evolution of variables of interest over
time}\label{evolution-of-variables-of-interest-over-time}

There are a number of values for the social class variables that are
outside the 1-7 scale. An interesting question is whether these
variables were re-coded or just dropped from the cleaned social mobility
dataset or multiply imputed entirely, and what these codes originally
mapped to.

\begin{Shaded}
\begin{Highlighting}[]
\CommentTok{\# I pretend that the "roughly ordinal" ordering of the social class variables}
\CommentTok{\# are strict orderings for visualization purposes}

\CommentTok{\# Filtering out suspected junk values for social class (outside the 1{-}7 scale)}
\NormalTok{wk\_data2 }\OtherTok{\textless{}{-}}\NormalTok{ wk\_data1 }\SpecialCharTok{\%\textgreater{}\%}
  \FunctionTok{mutate}\NormalTok{(}\AttributeTok{fclass =} \FunctionTok{as.numeric}\NormalTok{(fclass), }
         \AttributeTok{class38 =} \FunctionTok{as.numeric}\NormalTok{(class38)) }\SpecialCharTok{\%\textgreater{}\%}
  \FunctionTok{filter}\NormalTok{(class38 }\SpecialCharTok{\textgreater{}=} \DecValTok{1} \SpecialCharTok{\&}\NormalTok{ class38 }\SpecialCharTok{\textless{}=} \DecValTok{7}\NormalTok{,}
\NormalTok{         fclass }\SpecialCharTok{\textgreater{}=} \DecValTok{1} \SpecialCharTok{\&}\NormalTok{ class38 }\SpecialCharTok{\textless{}=} \DecValTok{7}\NormalTok{) }

\NormalTok{wk\_data2 }\SpecialCharTok{\%\textgreater{}\%}
  \FunctionTok{pivot\_longer}\NormalTok{(}\AttributeTok{cols =} \FunctionTok{c}\NormalTok{(ed37, class38, fclass), }\AttributeTok{names\_to =} \StringTok{"variable"}\NormalTok{, }\AttributeTok{values\_to =} \StringTok{"value"}\NormalTok{) }\SpecialCharTok{\%\textgreater{}\%}
  \FunctionTok{ggplot}\NormalTok{(}\FunctionTok{aes}\NormalTok{(}\AttributeTok{x =}\NormalTok{ value, }\AttributeTok{fill =}\NormalTok{ cohort)) }\SpecialCharTok{+}
  \FunctionTok{geom\_boxplot}\NormalTok{(}\AttributeTok{alpha =} \FloatTok{0.5}\NormalTok{) }\SpecialCharTok{+} 
  \FunctionTok{facet\_wrap}\NormalTok{(}\SpecialCharTok{\textasciitilde{}}\NormalTok{variable, }\AttributeTok{scales =} \StringTok{"free"}\NormalTok{) }\SpecialCharTok{+} 
  \FunctionTok{guides}\NormalTok{(}\AttributeTok{fill =} \FunctionTok{guide\_legend}\NormalTok{(}\AttributeTok{reverse =} \ConstantTok{TRUE}\NormalTok{)) }\SpecialCharTok{+} 
  \FunctionTok{theme}\NormalTok{(}\AttributeTok{axis.text.x =} \FunctionTok{element\_blank}\NormalTok{(),}
        \AttributeTok{axis.title.x =} \FunctionTok{element\_blank}\NormalTok{())}
\end{Highlighting}
\end{Shaded}

\pandocbounded{\includegraphics[keepaspectratio]{social_mobility_data_eda_files/figure-latex/unnamed-chunk-18-1.pdf}}

\begin{Shaded}
\begin{Highlighting}[]
\CommentTok{\# Not filtering out junk values for social class }
\NormalTok{wk\_data3 }\OtherTok{\textless{}{-}}\NormalTok{ wk\_data1 }\SpecialCharTok{\%\textgreater{}\%}
  \FunctionTok{mutate}\NormalTok{(}\AttributeTok{fclass =} \FunctionTok{as.numeric}\NormalTok{(fclass), }
         \AttributeTok{class38 =} \FunctionTok{as.numeric}\NormalTok{(class38)) }

\NormalTok{wk\_data3 }\SpecialCharTok{\%\textgreater{}\%}
  \FunctionTok{pivot\_longer}\NormalTok{(}\AttributeTok{cols =} \FunctionTok{c}\NormalTok{(ed37, class38, fclass), }\AttributeTok{names\_to =} \StringTok{"variable"}\NormalTok{, }\AttributeTok{values\_to =} \StringTok{"value"}\NormalTok{) }\SpecialCharTok{\%\textgreater{}\%}
  \FunctionTok{ggplot}\NormalTok{(}\FunctionTok{aes}\NormalTok{(}\AttributeTok{x =}\NormalTok{ value, }\AttributeTok{fill =}\NormalTok{ cohort)) }\SpecialCharTok{+}
  \FunctionTok{geom\_boxplot}\NormalTok{(}\AttributeTok{alpha =} \FloatTok{0.5}\NormalTok{) }\SpecialCharTok{+} 
  \FunctionTok{facet\_wrap}\NormalTok{(}\SpecialCharTok{\textasciitilde{}}\NormalTok{variable, }\AttributeTok{scales =} \StringTok{"free"}\NormalTok{) }\SpecialCharTok{+} 
  \FunctionTok{guides}\NormalTok{(}\AttributeTok{fill =} \FunctionTok{guide\_legend}\NormalTok{(}\AttributeTok{reverse =} \ConstantTok{TRUE}\NormalTok{)) }\SpecialCharTok{+} 
  \FunctionTok{theme}\NormalTok{(}\AttributeTok{axis.text.x =} \FunctionTok{element\_blank}\NormalTok{(),}
        \AttributeTok{axis.title.x =} \FunctionTok{element\_blank}\NormalTok{())}
\end{Highlighting}
\end{Shaded}

\begin{verbatim}
## Warning: Removed 3900 rows containing non-finite outside the scale range
## (`stat_boxplot()`).
\end{verbatim}

\pandocbounded{\includegraphics[keepaspectratio]{social_mobility_data_eda_files/figure-latex/unnamed-chunk-19-1.pdf}}

The box plots using the full data without filtering seems to be closer
to the cohort level distributions of these variables in the cleaned
social mobility data. This suggests that instead of being dropped, these
observations were somehow re-coded.

\subsection{External validity}\label{external-validity}

\subsubsection{Matching the work history and social mobility
datasets}\label{matching-the-work-history-and-social-mobility-datasets}

\begin{Shaded}
\begin{Highlighting}[]
\CommentTok{\# Trying to naively match on \textasciigrave{}.id\textasciigrave{} = \textasciigrave{}id\textasciigrave{} }
\NormalTok{matched1 }\OtherTok{\textless{}{-}}\NormalTok{ sm\_data1 }\SpecialCharTok{\%\textgreater{}\%} 
  \FunctionTok{mutate}\NormalTok{(}\AttributeTok{id =} \FunctionTok{as.integer}\NormalTok{(.id)) }\SpecialCharTok{\%\textgreater{}\%}
  \FunctionTok{select}\NormalTok{(}\SpecialCharTok{{-}}\NormalTok{.id) }\SpecialCharTok{\%\textgreater{}\%}
  \FunctionTok{left\_join}\NormalTok{(wk\_data1 }\SpecialCharTok{\%\textgreater{}\%} 
              \FunctionTok{filter}\NormalTok{(cohort }\SpecialCharTok{!=} \StringTok{"1970"}\NormalTok{), }\AttributeTok{by =} \FunctionTok{c}\NormalTok{(}\StringTok{"id"}\NormalTok{, }\StringTok{"cohort"}\NormalTok{))}
\CommentTok{\# Clear mismatches in the sex, social class and education variables}

\CommentTok{\# Trying to naively match on \textasciigrave{}id\textasciigrave{} = \textasciigrave{}id\textasciigrave{} }
\NormalTok{matched2 }\OtherTok{\textless{}{-}}\NormalTok{ sm\_data1 }\SpecialCharTok{\%\textgreater{}\%} 
  \FunctionTok{mutate}\NormalTok{(}\AttributeTok{id =} \FunctionTok{as.integer}\NormalTok{(id)) }\SpecialCharTok{\%\textgreater{}\%}
  \FunctionTok{select}\NormalTok{(}\SpecialCharTok{{-}}\NormalTok{.id) }\SpecialCharTok{\%\textgreater{}\%}
  \FunctionTok{left\_join}\NormalTok{(wk\_data1 }\SpecialCharTok{\%\textgreater{}\%} 
              \FunctionTok{filter}\NormalTok{(cohort }\SpecialCharTok{!=} \StringTok{"1970"}\NormalTok{), }\AttributeTok{by =} \FunctionTok{c}\NormalTok{(}\StringTok{"id"}\NormalTok{, }\StringTok{"cohort"}\NormalTok{))}
\CommentTok{\# Clear mismatches in the social class and education variables}
\end{Highlighting}
\end{Shaded}

A manual review of the datasets reveals that the correct matching scheme
seems to be match the row number of the stacked work history data with
the \texttt{.id} variable in the social mobility data.

\begin{Shaded}
\begin{Highlighting}[]
\NormalTok{wk\_full }\OtherTok{\textless{}{-}} \FunctionTok{bind\_rows}\NormalTok{(workhistory1 }\SpecialCharTok{\%\textgreater{}\%} \FunctionTok{mutate}\NormalTok{(}\AttributeTok{cohort =} \StringTok{"1946"}\NormalTok{,}
                                             \AttributeTok{ed37 =}\NormalTok{ educ),}
\NormalTok{                     workhistory2 }\SpecialCharTok{\%\textgreater{}\%} \FunctionTok{mutate}\NormalTok{(}\AttributeTok{cohort =} \StringTok{"1958"}\NormalTok{),}
\NormalTok{                     workhistory3 }\SpecialCharTok{\%\textgreater{}\%} \FunctionTok{mutate}\NormalTok{(}\AttributeTok{cohort =} \StringTok{"1970"}\NormalTok{)) }\SpecialCharTok{\%\textgreater{}\%}
  \FunctionTok{mutate}\NormalTok{(}\AttributeTok{id =} \FunctionTok{row\_number}\NormalTok{()) }\SpecialCharTok{\%\textgreater{}\%}
  \FunctionTok{select}\NormalTok{(cohort, id, sex, fclass, class38, ed37, IQ)}

\NormalTok{matched3 }\OtherTok{\textless{}{-}}\NormalTok{ sm\_data1 }\SpecialCharTok{\%\textgreater{}\%}
  \FunctionTok{mutate}\NormalTok{(}\AttributeTok{in\_sm =} \DecValTok{1}\NormalTok{,}
         \AttributeTok{id =}\NormalTok{ .id) }\SpecialCharTok{\%\textgreater{}\%}
  \FunctionTok{full\_join}\NormalTok{(wk\_full }\SpecialCharTok{\%\textgreater{}\%} \FunctionTok{mutate}\NormalTok{(}\AttributeTok{in\_wh =} \DecValTok{1}\NormalTok{,}
                               \AttributeTok{id =} \FunctionTok{as.factor}\NormalTok{(id)), }
            \AttributeTok{by =} \FunctionTok{c}\NormalTok{(}\StringTok{"id"}\NormalTok{, }\StringTok{"cohort"}\NormalTok{)) }\SpecialCharTok{\%\textgreater{}\%}
  \FunctionTok{mutate}\NormalTok{(}\AttributeTok{in\_both =} \FunctionTok{ifelse}\NormalTok{(}\SpecialCharTok{!}\FunctionTok{is.na}\NormalTok{(in\_sm) }\SpecialCharTok{\&} \SpecialCharTok{!}\FunctionTok{is.na}\NormalTok{(in\_wh), }\DecValTok{1}\NormalTok{, }\DecValTok{0}\NormalTok{),}
         \AttributeTok{sex.x =} \FunctionTok{ifelse}\NormalTok{(sex.x }\SpecialCharTok{==} \StringTok{"male"}\NormalTok{, }\DecValTok{0}\NormalTok{, }\DecValTok{1}\NormalTok{),}
         \AttributeTok{class38.x =} \FunctionTok{case\_when}\NormalTok{(class38.x }\SpecialCharTok{==} \StringTok{"High.sal"} \SpecialCharTok{\textasciitilde{}} \DecValTok{1}\NormalTok{,}
\NormalTok{                               class38.x }\SpecialCharTok{==} \StringTok{"Low.sal"} \SpecialCharTok{\textasciitilde{}} \DecValTok{2}\NormalTok{,}
\NormalTok{                               class38.x }\SpecialCharTok{==} \StringTok{"Interm"} \SpecialCharTok{\textasciitilde{}} \DecValTok{3}\NormalTok{,}
\NormalTok{                               class38.x }\SpecialCharTok{==} \StringTok{"Small.emp"} \SpecialCharTok{\textasciitilde{}} \DecValTok{4}\NormalTok{,}
\NormalTok{                               class38.x }\SpecialCharTok{==} \StringTok{"Low.sup"} \SpecialCharTok{\textasciitilde{}} \DecValTok{5}\NormalTok{,}
\NormalTok{                               class38.x }\SpecialCharTok{==} \StringTok{"Semi{-}rout"} \SpecialCharTok{\textasciitilde{}} \DecValTok{6}\NormalTok{,}
\NormalTok{                               class38.x }\SpecialCharTok{==} \StringTok{"Routine"} \SpecialCharTok{\textasciitilde{}} \DecValTok{7}\NormalTok{,}
\NormalTok{                               T }\SpecialCharTok{\textasciitilde{}} \ConstantTok{NA}\NormalTok{),}
         \AttributeTok{fclass.x =} \FunctionTok{case\_when}\NormalTok{(fclass.x }\SpecialCharTok{==} \StringTok{"High.sal"} \SpecialCharTok{\textasciitilde{}} \DecValTok{1}\NormalTok{,}
\NormalTok{                               fclass.x }\SpecialCharTok{==} \StringTok{"Low.sal"} \SpecialCharTok{\textasciitilde{}} \DecValTok{2}\NormalTok{,}
\NormalTok{                               fclass.x }\SpecialCharTok{==} \StringTok{"Interm"} \SpecialCharTok{\textasciitilde{}} \DecValTok{3}\NormalTok{,}
\NormalTok{                               fclass.x }\SpecialCharTok{==} \StringTok{"Small.emp"} \SpecialCharTok{\textasciitilde{}} \DecValTok{4}\NormalTok{,}
\NormalTok{                               fclass.x }\SpecialCharTok{==} \StringTok{"Low.sup"} \SpecialCharTok{\textasciitilde{}} \DecValTok{5}\NormalTok{,}
\NormalTok{                               fclass.x }\SpecialCharTok{==} \StringTok{"Semi{-}rout"} \SpecialCharTok{\textasciitilde{}} \DecValTok{6}\NormalTok{,}
\NormalTok{                               fclass.x }\SpecialCharTok{==} \StringTok{"Routine"} \SpecialCharTok{\textasciitilde{}} \DecValTok{7}\NormalTok{,}
\NormalTok{                               T }\SpecialCharTok{\textasciitilde{}} \ConstantTok{NA}\NormalTok{))}
\end{Highlighting}
\end{Shaded}

\begin{Shaded}
\begin{Highlighting}[]
\CommentTok{\# Checking agreement in auxiliary variables between the social mobility and work history datasets}

\NormalTok{check }\OtherTok{\textless{}{-}}\NormalTok{ matched3 }\SpecialCharTok{\%\textgreater{}\%}
  \FunctionTok{mutate}\NormalTok{(}\AttributeTok{fclass\_ag =}\NormalTok{ fclass.x }\SpecialCharTok{==}\NormalTok{ fclass.y,}
         \AttributeTok{educ37\_ag =}\NormalTok{ educ37 }\SpecialCharTok{==}\NormalTok{ ed37,}
         \AttributeTok{class38\_ag =}\NormalTok{ class38.x }\SpecialCharTok{==}\NormalTok{ class38.y)}

\FunctionTok{cat}\NormalTok{(}\StringTok{"Father\textquotesingle{}s class agreement: "}\NormalTok{, }\FunctionTok{round}\NormalTok{(}\FunctionTok{sum}\NormalTok{(check}\SpecialCharTok{$}\NormalTok{fclass\_ag, }\AttributeTok{na.rm =}\NormalTok{ T) }\SpecialCharTok{/} \FunctionTok{dim}\NormalTok{(check }\SpecialCharTok{\%\textgreater{}\%} \FunctionTok{drop\_na}\NormalTok{(fclass\_ag))[}\DecValTok{1}\NormalTok{], }\DecValTok{2}\NormalTok{), }\StringTok{"}\SpecialCharTok{\textbackslash{}n}\StringTok{"}\NormalTok{,}
    \StringTok{"Education agreement: "}\NormalTok{, }\FunctionTok{round}\NormalTok{(}\FunctionTok{sum}\NormalTok{(check}\SpecialCharTok{$}\NormalTok{educ37\_ag, }\AttributeTok{na.rm =}\NormalTok{ T) }\SpecialCharTok{/} \FunctionTok{dim}\NormalTok{(check }\SpecialCharTok{\%\textgreater{}\%} \FunctionTok{drop\_na}\NormalTok{(educ37\_ag))[}\DecValTok{1}\NormalTok{], }\DecValTok{2}\NormalTok{), }\StringTok{"}\SpecialCharTok{\textbackslash{}n}\StringTok{"}\NormalTok{,}
    \StringTok{"Unit class agreement: "}\NormalTok{, }\FunctionTok{round}\NormalTok{(}\FunctionTok{sum}\NormalTok{(check}\SpecialCharTok{$}\NormalTok{class38\_ag, }\AttributeTok{na.rm =}\NormalTok{ T) }\SpecialCharTok{/} \FunctionTok{dim}\NormalTok{(check }\SpecialCharTok{\%\textgreater{}\%} \FunctionTok{drop\_na}\NormalTok{(class38\_ag))[}\DecValTok{1}\NormalTok{], }\DecValTok{2}\NormalTok{), }\StringTok{"}\SpecialCharTok{\textbackslash{}n}\StringTok{"}\NormalTok{)}
\end{Highlighting}
\end{Shaded}

\begin{verbatim}
## Father's class agreement:  1 
##  Education agreement:  1 
##  Unit class agreement:  1
\end{verbatim}

\begin{Shaded}
\begin{Highlighting}[]
\CommentTok{\# Perfect agreement}
\end{Highlighting}
\end{Shaded}

\begin{Shaded}
\begin{Highlighting}[]
\NormalTok{matched }\OtherTok{\textless{}{-}}\NormalTok{ matched3 }\SpecialCharTok{\%\textgreater{}\%}
  \FunctionTok{select}\NormalTok{(id, cohort, weight, sex.y, fclass.y, class38.y, ed37, }
\NormalTok{         IQ, in\_wh, in\_both) }\SpecialCharTok{\%\textgreater{}\%}
  \FunctionTok{rename}\NormalTok{(}\AttributeTok{sex =}\NormalTok{ sex.y, }\AttributeTok{fclass =}\NormalTok{ fclass.y, }\AttributeTok{class =}\NormalTok{ class38.y,}
         \AttributeTok{educ =}\NormalTok{ ed37) }\SpecialCharTok{\%\textgreater{}\%}
  \FunctionTok{mutate}\NormalTok{(}\AttributeTok{class =} \FunctionTok{ifelse}\NormalTok{(class }\SpecialCharTok{\textgreater{}=} \DecValTok{1} \SpecialCharTok{\&}\NormalTok{ class }\SpecialCharTok{\textless{}=} \DecValTok{7}\NormalTok{, class, }\ConstantTok{NA}\NormalTok{))}

\CommentTok{\# There are observations in the work history 1946 data that are not in the social mobility data. This is because of the weighting that was performed to account for the sampling quirks of the 1946 data. For consistency, these observations will be dropped for now. }

\NormalTok{matched }\OtherTok{\textless{}{-}}\NormalTok{ matched }\SpecialCharTok{\%\textgreater{}\%} \FunctionTok{filter}\NormalTok{(in\_both }\SpecialCharTok{==} \DecValTok{1}\NormalTok{)}
\end{Highlighting}
\end{Shaded}

\subsubsection{Properties of the matched
dataset}\label{properties-of-the-matched-dataset}

\paragraph{Prevalence of missing
variables}\label{prevalence-of-missing-variables}

\begin{Shaded}
\begin{Highlighting}[]
\FunctionTok{skim}\NormalTok{(matched) }\SpecialCharTok{\%\textgreater{}\%}
  \FunctionTok{mutate}\NormalTok{(}\FunctionTok{across}\NormalTok{(}\FunctionTok{where}\NormalTok{(is.numeric), }\SpecialCharTok{\textasciitilde{}}\FunctionTok{round}\NormalTok{(., }\DecValTok{2}\NormalTok{)))}
\end{Highlighting}
\end{Shaded}

\begin{longtable}[]{@{}ll@{}}
\caption{Data summary}\tabularnewline
\toprule\noalign{}
\endfirsthead
\endhead
\bottomrule\noalign{}
\endlastfoot
Name & matched \\
Number of rows & 22579 \\
Number of columns & 10 \\
\_\_\_\_\_\_\_\_\_\_\_\_\_\_\_\_\_\_\_\_\_\_\_ & \\
Column type frequency: & \\
character & 1 \\
factor & 1 \\
numeric & 8 \\
\_\_\_\_\_\_\_\_\_\_\_\_\_\_\_\_\_\_\_\_\_\_\_\_ & \\
Group variables & None \\
\end{longtable}

\textbf{Variable type: character}

\begin{longtable}[]{@{}
  >{\raggedright\arraybackslash}p{(\linewidth - 14\tabcolsep) * \real{0.1944}}
  >{\raggedleft\arraybackslash}p{(\linewidth - 14\tabcolsep) * \real{0.1389}}
  >{\raggedleft\arraybackslash}p{(\linewidth - 14\tabcolsep) * \real{0.1944}}
  >{\raggedleft\arraybackslash}p{(\linewidth - 14\tabcolsep) * \real{0.0556}}
  >{\raggedleft\arraybackslash}p{(\linewidth - 14\tabcolsep) * \real{0.0556}}
  >{\raggedleft\arraybackslash}p{(\linewidth - 14\tabcolsep) * \real{0.0833}}
  >{\raggedleft\arraybackslash}p{(\linewidth - 14\tabcolsep) * \real{0.1250}}
  >{\raggedleft\arraybackslash}p{(\linewidth - 14\tabcolsep) * \real{0.1528}}@{}}
\toprule\noalign{}
\begin{minipage}[b]{\linewidth}\raggedright
skim\_variable
\end{minipage} & \begin{minipage}[b]{\linewidth}\raggedleft
n\_missing
\end{minipage} & \begin{minipage}[b]{\linewidth}\raggedleft
complete\_rate
\end{minipage} & \begin{minipage}[b]{\linewidth}\raggedleft
min
\end{minipage} & \begin{minipage}[b]{\linewidth}\raggedleft
max
\end{minipage} & \begin{minipage}[b]{\linewidth}\raggedleft
empty
\end{minipage} & \begin{minipage}[b]{\linewidth}\raggedleft
n\_unique
\end{minipage} & \begin{minipage}[b]{\linewidth}\raggedleft
whitespace
\end{minipage} \\
\midrule\noalign{}
\endhead
\bottomrule\noalign{}
\endlastfoot
cohort & 0 & 1 & 4 & 4 & 0 & 3 & 0 \\
\end{longtable}

\textbf{Variable type: factor}

\begin{longtable}[]{@{}
  >{\raggedright\arraybackslash}p{(\linewidth - 10\tabcolsep) * \real{0.1647}}
  >{\raggedleft\arraybackslash}p{(\linewidth - 10\tabcolsep) * \real{0.1176}}
  >{\raggedleft\arraybackslash}p{(\linewidth - 10\tabcolsep) * \real{0.1647}}
  >{\raggedright\arraybackslash}p{(\linewidth - 10\tabcolsep) * \real{0.0941}}
  >{\raggedleft\arraybackslash}p{(\linewidth - 10\tabcolsep) * \real{0.1059}}
  >{\raggedright\arraybackslash}p{(\linewidth - 10\tabcolsep) * \real{0.3529}}@{}}
\toprule\noalign{}
\begin{minipage}[b]{\linewidth}\raggedright
skim\_variable
\end{minipage} & \begin{minipage}[b]{\linewidth}\raggedleft
n\_missing
\end{minipage} & \begin{minipage}[b]{\linewidth}\raggedleft
complete\_rate
\end{minipage} & \begin{minipage}[b]{\linewidth}\raggedright
ordered
\end{minipage} & \begin{minipage}[b]{\linewidth}\raggedleft
n\_unique
\end{minipage} & \begin{minipage}[b]{\linewidth}\raggedright
top\_counts
\end{minipage} \\
\midrule\noalign{}
\endhead
\bottomrule\noalign{}
\endlastfoot
id & 0 & 1 & FALSE & 22579 & 10: 1, 100: 1, 100: 1, 100: 1 \\
\end{longtable}

\textbf{Variable type: numeric}

\begin{longtable}[]{@{}
  >{\raggedright\arraybackslash}p{(\linewidth - 20\tabcolsep) * \real{0.1728}}
  >{\raggedleft\arraybackslash}p{(\linewidth - 20\tabcolsep) * \real{0.1235}}
  >{\raggedleft\arraybackslash}p{(\linewidth - 20\tabcolsep) * \real{0.1728}}
  >{\raggedleft\arraybackslash}p{(\linewidth - 20\tabcolsep) * \real{0.0617}}
  >{\raggedleft\arraybackslash}p{(\linewidth - 20\tabcolsep) * \real{0.0617}}
  >{\raggedleft\arraybackslash}p{(\linewidth - 20\tabcolsep) * \real{0.0741}}
  >{\raggedleft\arraybackslash}p{(\linewidth - 20\tabcolsep) * \real{0.0741}}
  >{\raggedleft\arraybackslash}p{(\linewidth - 20\tabcolsep) * \real{0.0617}}
  >{\raggedleft\arraybackslash}p{(\linewidth - 20\tabcolsep) * \real{0.0617}}
  >{\raggedleft\arraybackslash}p{(\linewidth - 20\tabcolsep) * \real{0.0617}}
  >{\raggedright\arraybackslash}p{(\linewidth - 20\tabcolsep) * \real{0.0741}}@{}}
\toprule\noalign{}
\begin{minipage}[b]{\linewidth}\raggedright
skim\_variable
\end{minipage} & \begin{minipage}[b]{\linewidth}\raggedleft
n\_missing
\end{minipage} & \begin{minipage}[b]{\linewidth}\raggedleft
complete\_rate
\end{minipage} & \begin{minipage}[b]{\linewidth}\raggedleft
mean
\end{minipage} & \begin{minipage}[b]{\linewidth}\raggedleft
sd
\end{minipage} & \begin{minipage}[b]{\linewidth}\raggedleft
p0
\end{minipage} & \begin{minipage}[b]{\linewidth}\raggedleft
p25
\end{minipage} & \begin{minipage}[b]{\linewidth}\raggedleft
p50
\end{minipage} & \begin{minipage}[b]{\linewidth}\raggedleft
p75
\end{minipage} & \begin{minipage}[b]{\linewidth}\raggedleft
p100
\end{minipage} & \begin{minipage}[b]{\linewidth}\raggedright
hist
\end{minipage} \\
\midrule\noalign{}
\endhead
\bottomrule\noalign{}
\endlastfoot
weight & 0 & 1.00 & 1.00 & 0.26 & 0.45 & 1.00 & 1.00 & 1.00 & 1.80 &
▁▁▇▁▁ \\
sex & 0 & 1.00 & 0.31 & 0.46 & 0.00 & 0.00 & 0.00 & 1.00 & 1.00 &
▇▁▁▁▃ \\
fclass & 2812 & 0.88 & 4.46 & 2.07 & 1.00 & 2.00 & 5.00 & 7.00 & 7.00 &
▅▂▂▃▇ \\
class & 7099 & 0.69 & 3.57 & 2.02 & 1.00 & 2.00 & 3.00 & 6.00 & 7.00 &
▇▃▂▂▅ \\
educ & 84 & 1.00 & 3.78 & 2.12 & 1.00 & 2.00 & 4.00 & 6.00 & 8.00 &
▇▆▆▃▅ \\
IQ & 7582 & 0.66 & 0.03 & 1.00 & -3.66 & -0.68 & 0.07 & 0.77 & 3.28 &
▁▃▇▆▁ \\
in\_wh & 0 & 1.00 & 1.00 & 0.00 & 1.00 & 1.00 & 1.00 & 1.00 & 1.00 &
▁▁▇▁▁ \\
in\_both & 0 & 1.00 & 1.00 & 0.00 & 1.00 & 1.00 & 1.00 & 1.00 & 1.00 &
▁▁▇▁▁ \\
\end{longtable}

\subparagraph{Checking the distribution of other variables based on
whether they are missing the IQ variable or
not}\label{checking-the-distribution-of-other-variables-based-on-whether-they-are-missing-the-iq-variable-or-not}

\begin{Shaded}
\begin{Highlighting}[]
\FunctionTok{summary}\NormalTok{(matched }\SpecialCharTok{\%\textgreater{}\%} \FunctionTok{filter}\NormalTok{(}\FunctionTok{is.na}\NormalTok{(IQ)))}
\end{Highlighting}
\end{Shaded}

\begin{verbatim}
##        id          cohort              weight            sex        
##  10003  :   1   Length:7582        Min.   :0.4512   Min.   :0.0000  
##  10004  :   1   Class :character   1st Qu.:1.0000   1st Qu.:0.0000  
##  1001   :   1   Mode  :character   Median :1.0000   Median :0.0000  
##  10014  :   1                      Mean   :0.9930   Mean   :0.3168  
##  10016  :   1                      3rd Qu.:1.0000   3rd Qu.:1.0000  
##  10017  :   1                      Max.   :1.8047   Max.   :1.0000  
##  (Other):7576                                                       
##      fclass          class            educ             IQ           in_wh  
##  Min.   :1.000   Min.   :1.000   Min.   :1.000   Min.   : NA    Min.   :1  
##  1st Qu.:2.000   1st Qu.:2.000   1st Qu.:2.000   1st Qu.: NA    1st Qu.:1  
##  Median :5.000   Median :3.000   Median :3.000   Median : NA    Median :1  
##  Mean   :4.552   Mean   :3.635   Mean   :3.684   Mean   :NaN    Mean   :1  
##  3rd Qu.:7.000   3rd Qu.:6.000   3rd Qu.:6.000   3rd Qu.: NA    3rd Qu.:1  
##  Max.   :7.000   Max.   :7.000   Max.   :8.000   Max.   : NA    Max.   :1  
##  NA's   :1902    NA's   :2429    NA's   :77      NA's   :7582              
##     in_both 
##  Min.   :1  
##  1st Qu.:1  
##  Median :1  
##  Mean   :1  
##  3rd Qu.:1  
##  Max.   :1  
## 
\end{verbatim}

\begin{Shaded}
\begin{Highlighting}[]
\FunctionTok{summary}\NormalTok{(matched }\SpecialCharTok{\%\textgreater{}\%} \FunctionTok{filter}\NormalTok{(}\SpecialCharTok{!}\FunctionTok{is.na}\NormalTok{(IQ)))}
\end{Highlighting}
\end{Shaded}

\begin{verbatim}
##        id           cohort              weight            sex        
##  10     :    1   Length:14997       Min.   :0.4512   Min.   :0.0000  
##  100    :    1   Class :character   1st Qu.:1.0000   1st Qu.:0.0000  
##  10000  :    1   Mode  :character   Median :1.0000   Median :0.0000  
##  10005  :    1                      Mean   :0.9990   Mean   :0.3057  
##  10006  :    1                      3rd Qu.:1.0000   3rd Qu.:1.0000  
##  10008  :    1                      Max.   :1.8047   Max.   :1.0000  
##  (Other):14991                                                       
##      fclass          class            educ             IQ               in_wh  
##  Min.   :1.000   Min.   :1.000   Min.   :1.000   Min.   :-3.65914   Min.   :1  
##  1st Qu.:3.000   1st Qu.:2.000   1st Qu.:2.000   1st Qu.:-0.68476   1st Qu.:1  
##  Median :5.000   Median :3.000   Median :4.000   Median : 0.07289   Median :1  
##  Mean   :4.429   Mean   :3.537   Mean   :3.833   Mean   : 0.02705   Mean   :1  
##  3rd Qu.:6.000   3rd Qu.:5.000   3rd Qu.:6.000   3rd Qu.: 0.77292   3rd Qu.:1  
##  Max.   :7.000   Max.   :7.000   Max.   :8.000   Max.   : 3.28085   Max.   :1  
##  NA's   :910     NA's   :4670    NA's   :7                                     
##     in_both 
##  Min.   :1  
##  1st Qu.:1  
##  Median :1  
##  Mean   :1  
##  3rd Qu.:1  
##  Max.   :1  
## 
\end{verbatim}

No significant differences in other variables observed

\subparagraph{Checking the distribution of other variables based on
whether they are missing the unit class variable or
not}\label{checking-the-distribution-of-other-variables-based-on-whether-they-are-missing-the-unit-class-variable-or-not}

\begin{Shaded}
\begin{Highlighting}[]
\FunctionTok{summary}\NormalTok{(matched }\SpecialCharTok{\%\textgreater{}\%} \FunctionTok{filter}\NormalTok{(}\FunctionTok{is.na}\NormalTok{(class)))}
\end{Highlighting}
\end{Shaded}

\begin{verbatim}
##        id          cohort              weight            sex        
##  100    :   1   Length:7099        Min.   :0.4512   Min.   :0.0000  
##  1001   :   1   Class :character   1st Qu.:1.0000   1st Qu.:0.0000  
##  10011  :   1   Mode  :character   Median :1.0000   Median :0.0000  
##  10017  :   1                      Mean   :0.9811   Mean   :0.3364  
##  10021  :   1                      3rd Qu.:1.0000   3rd Qu.:1.0000  
##  10024  :   1                      Max.   :1.8047   Max.   :1.0000  
##  (Other):7093                                                       
##      fclass          class           educ             IQ              in_wh  
##  Min.   :1.000   Min.   : NA    Min.   :1.000   Min.   :-3.6591   Min.   :1  
##  1st Qu.:3.000   1st Qu.: NA    1st Qu.:1.000   1st Qu.:-0.9215   1st Qu.:1  
##  Median :5.000   Median : NA    Median :3.000   Median :-0.1419   Median :1  
##  Mean   :4.687   Mean   :NaN    Mean   :3.318   Mean   :-0.1712   Mean   :1  
##  3rd Qu.:7.000   3rd Qu.: NA    3rd Qu.:4.000   3rd Qu.: 0.6026   3rd Qu.:1  
##  Max.   :7.000   Max.   : NA    Max.   :8.000   Max.   : 3.2555   Max.   :1  
##  NA's   :1131    NA's   :7099   NA's   :52      NA's   :2429                 
##     in_both 
##  Min.   :1  
##  1st Qu.:1  
##  Median :1  
##  Mean   :1  
##  3rd Qu.:1  
##  Max.   :1  
## 
\end{verbatim}

\begin{Shaded}
\begin{Highlighting}[]
\FunctionTok{summary}\NormalTok{(matched }\SpecialCharTok{\%\textgreater{}\%} \FunctionTok{filter}\NormalTok{(}\SpecialCharTok{!}\FunctionTok{is.na}\NormalTok{(class)))}
\end{Highlighting}
\end{Shaded}

\begin{verbatim}
##        id           cohort              weight            sex        
##  10     :    1   Length:15480       Min.   :0.4512   Min.   :0.0000  
##  10000  :    1   Class :character   1st Qu.:1.0000   1st Qu.:0.0000  
##  10003  :    1   Mode  :character   Median :1.0000   Median :0.0000  
##  10004  :    1                      Mean   :1.0043   Mean   :0.2971  
##  10005  :    1                      3rd Qu.:1.0000   3rd Qu.:1.0000  
##  10006  :    1                      Max.   :1.8047   Max.   :1.0000  
##  (Other):15474                                                       
##      fclass          class           educ             IQ              in_wh  
##  Min.   :1.000   Min.   :1.00   Min.   :1.000   Min.   :-3.2823   Min.   :1  
##  1st Qu.:2.000   1st Qu.:2.00   1st Qu.:2.000   1st Qu.:-0.5685   1st Qu.:1  
##  Median :5.000   Median :3.00   Median :4.000   Median : 0.1617   Median :1  
##  Mean   :4.368   Mean   :3.57   Mean   :3.996   Mean   : 0.1167   Mean   :1  
##  3rd Qu.:6.000   3rd Qu.:6.00   3rd Qu.:6.000   3rd Qu.: 0.8290   3rd Qu.:1  
##  Max.   :7.000   Max.   :7.00   Max.   :8.000   Max.   : 3.2808   Max.   :1  
##  NA's   :1681                   NA's   :32      NA's   :5153                 
##     in_both 
##  Min.   :1  
##  1st Qu.:1  
##  Median :1  
##  Mean   :1  
##  3rd Qu.:1  
##  Max.   :1  
## 
\end{verbatim}

Larger differences are observed here, where is unit's social class is
not observed, it seems that these units are in ``lower'' social classes
and have lower average levels of education.

\subparagraph{Dropping variables with any missing
variables}\label{dropping-variables-with-any-missing-variables}

The data with missing variables can be handled by missing imputation
procedures if this becomes a focus of my dissertation. However, because
I am mostly using this dataset for illustrative purposes, I will work
with the ``complete case analysis'' data for now.

\begin{Shaded}
\begin{Highlighting}[]
\NormalTok{data\_cca }\OtherTok{\textless{}{-}}\NormalTok{ matched }\SpecialCharTok{\%\textgreater{}\%} 
  \FunctionTok{drop\_na}\NormalTok{() }\SpecialCharTok{\%\textgreater{}\%}
  \FunctionTok{select}\NormalTok{(}\SpecialCharTok{{-}}\NormalTok{in\_wh, }\SpecialCharTok{{-}}\NormalTok{in\_both) }\SpecialCharTok{\%\textgreater{}\%}
  \FunctionTok{mutate}\NormalTok{(}\AttributeTok{cohort =} \FunctionTok{as.numeric}\NormalTok{(cohort))}
\end{Highlighting}
\end{Shaded}

\begin{Shaded}
\begin{Highlighting}[]
\FunctionTok{skim}\NormalTok{(data\_cca) }\SpecialCharTok{\%\textgreater{}\%}
  \FunctionTok{mutate}\NormalTok{(}\FunctionTok{across}\NormalTok{(}\FunctionTok{where}\NormalTok{(is.numeric), }\SpecialCharTok{\textasciitilde{}}\FunctionTok{round}\NormalTok{(., }\DecValTok{2}\NormalTok{)))}
\end{Highlighting}
\end{Shaded}

\begin{longtable}[]{@{}ll@{}}
\caption{Data summary}\tabularnewline
\toprule\noalign{}
\endfirsthead
\endhead
\bottomrule\noalign{}
\endlastfoot
Name & data\_cca \\
Number of rows & 9803 \\
Number of columns & 8 \\
\_\_\_\_\_\_\_\_\_\_\_\_\_\_\_\_\_\_\_\_\_\_\_ & \\
Column type frequency: & \\
factor & 1 \\
numeric & 7 \\
\_\_\_\_\_\_\_\_\_\_\_\_\_\_\_\_\_\_\_\_\_\_\_\_ & \\
Group variables & None \\
\end{longtable}

\textbf{Variable type: factor}

\begin{longtable}[]{@{}
  >{\raggedright\arraybackslash}p{(\linewidth - 10\tabcolsep) * \real{0.1647}}
  >{\raggedleft\arraybackslash}p{(\linewidth - 10\tabcolsep) * \real{0.1176}}
  >{\raggedleft\arraybackslash}p{(\linewidth - 10\tabcolsep) * \real{0.1647}}
  >{\raggedright\arraybackslash}p{(\linewidth - 10\tabcolsep) * \real{0.0941}}
  >{\raggedleft\arraybackslash}p{(\linewidth - 10\tabcolsep) * \real{0.1059}}
  >{\raggedright\arraybackslash}p{(\linewidth - 10\tabcolsep) * \real{0.3529}}@{}}
\toprule\noalign{}
\begin{minipage}[b]{\linewidth}\raggedright
skim\_variable
\end{minipage} & \begin{minipage}[b]{\linewidth}\raggedleft
n\_missing
\end{minipage} & \begin{minipage}[b]{\linewidth}\raggedleft
complete\_rate
\end{minipage} & \begin{minipage}[b]{\linewidth}\raggedright
ordered
\end{minipage} & \begin{minipage}[b]{\linewidth}\raggedleft
n\_unique
\end{minipage} & \begin{minipage}[b]{\linewidth}\raggedright
top\_counts
\end{minipage} \\
\midrule\noalign{}
\endhead
\bottomrule\noalign{}
\endlastfoot
id & 0 & 1 & FALSE & 9803 & 10: 1, 100: 1, 100: 1, 100: 1 \\
\end{longtable}

\textbf{Variable type: numeric}

\begin{longtable}[]{@{}
  >{\raggedright\arraybackslash}p{(\linewidth - 20\tabcolsep) * \real{0.1443}}
  >{\raggedleft\arraybackslash}p{(\linewidth - 20\tabcolsep) * \real{0.1031}}
  >{\raggedleft\arraybackslash}p{(\linewidth - 20\tabcolsep) * \real{0.1443}}
  >{\raggedleft\arraybackslash}p{(\linewidth - 20\tabcolsep) * \real{0.0825}}
  >{\raggedleft\arraybackslash}p{(\linewidth - 20\tabcolsep) * \real{0.0515}}
  >{\raggedleft\arraybackslash}p{(\linewidth - 20\tabcolsep) * \real{0.0825}}
  >{\raggedleft\arraybackslash}p{(\linewidth - 20\tabcolsep) * \real{0.0825}}
  >{\raggedleft\arraybackslash}p{(\linewidth - 20\tabcolsep) * \real{0.0825}}
  >{\raggedleft\arraybackslash}p{(\linewidth - 20\tabcolsep) * \real{0.0825}}
  >{\raggedleft\arraybackslash}p{(\linewidth - 20\tabcolsep) * \real{0.0825}}
  >{\raggedright\arraybackslash}p{(\linewidth - 20\tabcolsep) * \real{0.0619}}@{}}
\toprule\noalign{}
\begin{minipage}[b]{\linewidth}\raggedright
skim\_variable
\end{minipage} & \begin{minipage}[b]{\linewidth}\raggedleft
n\_missing
\end{minipage} & \begin{minipage}[b]{\linewidth}\raggedleft
complete\_rate
\end{minipage} & \begin{minipage}[b]{\linewidth}\raggedleft
mean
\end{minipage} & \begin{minipage}[b]{\linewidth}\raggedleft
sd
\end{minipage} & \begin{minipage}[b]{\linewidth}\raggedleft
p0
\end{minipage} & \begin{minipage}[b]{\linewidth}\raggedleft
p25
\end{minipage} & \begin{minipage}[b]{\linewidth}\raggedleft
p50
\end{minipage} & \begin{minipage}[b]{\linewidth}\raggedleft
p75
\end{minipage} & \begin{minipage}[b]{\linewidth}\raggedleft
p100
\end{minipage} & \begin{minipage}[b]{\linewidth}\raggedright
hist
\end{minipage} \\
\midrule\noalign{}
\endhead
\bottomrule\noalign{}
\endlastfoot
cohort & 0 & 1 & 1960.89 & 8.36 & 1946.00 & 1958.00 & 1958.00 & 1970.00
& 1970.00 & ▂▁▇▁▇ \\
weight & 0 & 1 & 1.01 & 0.26 & 0.45 & 1.00 & 1.00 & 1.00 & 1.80 &
▁▁▇▁▁ \\
sex & 0 & 1 & 0.30 & 0.46 & 0.00 & 0.00 & 0.00 & 1.00 & 1.00 & ▇▁▁▁▃ \\
fclass & 0 & 1 & 4.34 & 1.97 & 1.00 & 2.00 & 5.00 & 6.00 & 7.00 &
▆▃▂▅▇ \\
class & 0 & 1 & 3.53 & 2.01 & 1.00 & 2.00 & 3.00 & 5.00 & 7.00 &
▇▃▂▂▅ \\
educ & 0 & 1 & 4.06 & 2.08 & 1.00 & 3.00 & 4.00 & 6.00 & 8.00 & ▇▆▇▅▆ \\
IQ & 0 & 1 & 0.13 & 0.97 & -3.28 & -0.56 & 0.17 & 0.84 & 3.28 & ▁▃▇▅▁ \\
\end{longtable}

\begin{Shaded}
\begin{Highlighting}[]
\CommentTok{\# Number of observations broken by cohort and sex }
\NormalTok{data\_cca }\SpecialCharTok{\%\textgreater{}\%}
  \FunctionTok{group\_by}\NormalTok{(cohort, sex) }\SpecialCharTok{\%\textgreater{}\%}
  \FunctionTok{summarize}\NormalTok{(}\AttributeTok{count =} \FunctionTok{n}\NormalTok{()) }\SpecialCharTok{\%\textgreater{}\%}
  \FunctionTok{pivot\_wider}\NormalTok{(}\AttributeTok{names\_from =}\NormalTok{ cohort, }\AttributeTok{values\_from =}\NormalTok{ count)}
\end{Highlighting}
\end{Shaded}

\begin{verbatim}
## `summarise()` has regrouped the output.
## i Summaries were computed grouped by cohort and sex.
## i Output is grouped by cohort.
## i Use `summarise(.groups = "drop_last")` to silence this message.
## i Use `summarise(.by = c(cohort, sex))` for per-operation grouping
##   (`?dplyr::dplyr_by`) instead.
\end{verbatim}

\begin{verbatim}
## # A tibble: 2 x 4
##   sex        `1946` `1958` `1970`
##   <dbl+lbl>   <int>  <int>  <int>
## 1 0 [male]      969   3188   2730
## 2 1 [female]    513   1288   1115
\end{verbatim}

\paragraph{Checks for overlap of the IQ and education
variables}\label{checks-for-overlap-of-the-iq-and-education-variables}

The differential distributions of IQ and education by social classes is
of particular interest for my application of mediation analysis
approaches for analyzing social class driven disparities. This is
because for interventional effect approaches to mediation analysis, a
key assumption for estimators to be consistent is an analogue of the
common support assumption often adopted in causal inference selection on
observables. If the distribution of these two variables diverges too
much across (father's) social classes, certain non-parametric methods of
estimating interventional effects, like the ones advances by Benkesar
and Ran (2021) and Diaz et al (2021), become much less attractive and
more parametric assumptions need to be made to estimate the effects.

I expect that because education is a discrete variable, there will be
satisfactory overlap here. For the continuous IQ variable, however,
there will be more of a problem with overlap for estimation. Since the
goal of my thesis is to develop the theory (identification, estimation)
of a certain class of interventional effects, the contrast between these
two variables is actually very helpful for illustration of some of the
issues with common support that I will be discussing, and can justify
the choice of estimators I develop.

I will produce distributions for education and IQ by father's social
class both as a whole and broken down by the sex and cohort covariates.
Having overlap conditional on covariates that might confound the
relationship between IQ/education and unit social class (confounders of
the mediator-outcome relationship) is essential for identification, so I
will be interested in both the overall distributions and the cohort/sex
specific ones.

\begin{Shaded}
\begin{Highlighting}[]
\NormalTok{data\_cca }\OtherTok{\textless{}{-}}\NormalTok{ data\_cca }\SpecialCharTok{\%\textgreater{}\%} 
  \FunctionTok{mutate}\NormalTok{(}\AttributeTok{class\_c =} \FunctionTok{case\_when}\NormalTok{(class }\SpecialCharTok{==} \DecValTok{1} \SpecialCharTok{\textasciitilde{}} \StringTok{"High.sal"}\NormalTok{,}
\NormalTok{                               class }\SpecialCharTok{==} \DecValTok{2} \SpecialCharTok{\textasciitilde{}} \StringTok{"Low.sal"}\NormalTok{,}
\NormalTok{                               class }\SpecialCharTok{==} \DecValTok{3} \SpecialCharTok{\textasciitilde{}} \StringTok{"Interm"}\NormalTok{,}
\NormalTok{                               class }\SpecialCharTok{==} \DecValTok{4} \SpecialCharTok{\textasciitilde{}} \StringTok{"Small.emp"}\NormalTok{,}
\NormalTok{                               class }\SpecialCharTok{==} \DecValTok{5} \SpecialCharTok{\textasciitilde{}} \StringTok{"Low.sup"}\NormalTok{,}
\NormalTok{                               class }\SpecialCharTok{==} \DecValTok{6} \SpecialCharTok{\textasciitilde{}}\StringTok{"Semi{-}rout"}\NormalTok{,}
\NormalTok{                               class }\SpecialCharTok{==} \DecValTok{7} \SpecialCharTok{\textasciitilde{}} \StringTok{"Routine"}\NormalTok{,}
\NormalTok{                               T }\SpecialCharTok{\textasciitilde{}} \ConstantTok{NA}\NormalTok{),}
         \AttributeTok{fclass\_c =} \FunctionTok{case\_when}\NormalTok{(fclass }\SpecialCharTok{==} \DecValTok{1} \SpecialCharTok{\textasciitilde{}} \StringTok{"High.sal"}\NormalTok{,}
\NormalTok{                               fclass }\SpecialCharTok{==} \DecValTok{2} \SpecialCharTok{\textasciitilde{}} \StringTok{"Low.sal"}\NormalTok{,}
\NormalTok{                               fclass }\SpecialCharTok{==} \DecValTok{3} \SpecialCharTok{\textasciitilde{}} \StringTok{"Interm"}\NormalTok{,}
\NormalTok{                               fclass }\SpecialCharTok{==} \DecValTok{4} \SpecialCharTok{\textasciitilde{}} \StringTok{"Small.emp"}\NormalTok{,}
\NormalTok{                               fclass }\SpecialCharTok{==} \DecValTok{5} \SpecialCharTok{\textasciitilde{}} \StringTok{"Low.sup"}\NormalTok{,}
\NormalTok{                               fclass }\SpecialCharTok{==} \DecValTok{6} \SpecialCharTok{\textasciitilde{}}\StringTok{"Semi{-}rout"}\NormalTok{,}
\NormalTok{                               fclass }\SpecialCharTok{==} \DecValTok{7} \SpecialCharTok{\textasciitilde{}} \StringTok{"Routine"}\NormalTok{,}
\NormalTok{                               T }\SpecialCharTok{\textasciitilde{}} \ConstantTok{NA}\NormalTok{))}
\end{Highlighting}
\end{Shaded}

\subparagraph{Education distributions}\label{education-distributions}

\begin{Shaded}
\begin{Highlighting}[]
\NormalTok{overlap\_table}
\end{Highlighting}
\end{Shaded}

\begin{verbatim}
##  cohort    sex  1  2  3   4  5   6  7  8
##    1946   Male  2  3  2  13  2  12  4  1
##    1946 Female  2  5  4   9  1   2  1  1
##    1958   Male  7  4 16  23  4  23  9  2
##    1958 Female  5  4  6  10  1   9  4  1
##    1970   Male 19  8 29  51  6  46 39  8
##    1970 Female  3  4 10  19  1  15 13  2
##   Total      - 38 28 67 125 15 107 70 15
\end{verbatim}

There is complete overlap in education across all father social classes
even when broken down by cohort and sex. However, there will probably be
potential issues with the variance of inverse probability weighting
(IPW) based methods to estimate interventional effects conditional on
confounders because of the fact that some father's social classes will
have only 1 or 2 observations for a given level of education to based
estimates off of.

\subparagraph{IQ distributions}\label{iq-distributions}

Because this is a continuous variable, clearly there will be non-overlap
in the strict sense. The two options are to either assume that the
distribution of the IQ tests follows a parametrically specified
distribution in each group or to perform discretization on the
continuous data to satisfy positivity. The smaller the differences in
performance on the test, the more justified a discetization +
non-parametric approach would be. However, it is well known in the
literature that performance on IQ tests should be normally distributed
in the population (by construction), so a parametric assumption on the
structure of the distribution we are trying to estimate seems justified
here.

Here, I think it is helpful to first get a sense of the differences in
the distributions of IQ test scores by father's social class by looking
at some histograms.

\begin{Shaded}
\begin{Highlighting}[]
\FunctionTok{ggplot}\NormalTok{(}\AttributeTok{data =}\NormalTok{ data\_cca, }\FunctionTok{aes}\NormalTok{(}\AttributeTok{x =}\NormalTok{ IQ, }\AttributeTok{color =}\NormalTok{ fclass\_c)) }\SpecialCharTok{+}
  \FunctionTok{geom\_density}\NormalTok{(}\AttributeTok{alpha =} \FloatTok{0.5}\NormalTok{) }\SpecialCharTok{+}
  \FunctionTok{theme\_classic}\NormalTok{() }\SpecialCharTok{+}
  \FunctionTok{labs}\NormalTok{(}
    \AttributeTok{title =} \StringTok{""}\NormalTok{,}
    \AttributeTok{x =} \StringTok{"IQ test performance"}\NormalTok{,}
    \AttributeTok{y =} \StringTok{"Density"}
\NormalTok{  )}
\end{Highlighting}
\end{Shaded}

\pandocbounded{\includegraphics[keepaspectratio]{social_mobility_data_eda_files/figure-latex/unnamed-chunk-35-1.pdf}}

\begin{Shaded}
\begin{Highlighting}[]
\FunctionTok{ggplot}\NormalTok{(}\AttributeTok{data =}\NormalTok{ data\_cca }\SpecialCharTok{\%\textgreater{}\%} 
         \FunctionTok{mutate}\NormalTok{(}\AttributeTok{sex =} \FunctionTok{ifelse}\NormalTok{(sex }\SpecialCharTok{==} \DecValTok{0}\NormalTok{, }\StringTok{"Male"}\NormalTok{, }\StringTok{"Female"}\NormalTok{)), }
       \FunctionTok{aes}\NormalTok{(}\AttributeTok{x =}\NormalTok{ IQ, }\AttributeTok{color =}\NormalTok{ fclass\_c)) }\SpecialCharTok{+}
  \FunctionTok{geom\_density}\NormalTok{(}\AttributeTok{alpha =} \FloatTok{0.5}\NormalTok{) }\SpecialCharTok{+}
  \FunctionTok{theme\_classic}\NormalTok{() }\SpecialCharTok{+}
  \FunctionTok{labs}\NormalTok{(}
    \AttributeTok{title =} \StringTok{""}\NormalTok{,}
    \AttributeTok{x =} \StringTok{"IQ test performance"}\NormalTok{,}
    \AttributeTok{y =} \StringTok{"Density"}
\NormalTok{  ) }\SpecialCharTok{+}
  \FunctionTok{facet\_grid}\NormalTok{(sex}\SpecialCharTok{\textasciitilde{}}\NormalTok{cohort, }\AttributeTok{scales =} \StringTok{"free"}\NormalTok{)}
\end{Highlighting}
\end{Shaded}

\pandocbounded{\includegraphics[keepaspectratio]{social_mobility_data_eda_files/figure-latex/unnamed-chunk-36-1.pdf}}

Interestingly, there are larger observed discrepancies in IQ test
performance for the 1946 and 1958 cohorts compared to 1970, which looks
like it follows the relatively normal, minor discrepancies across
father's social class of the overall data. This might point to the
smaller sample size in 1946, but it also might be evidence of a change
in how the IQ test was administered and what it consisted of across the
24 years of the data. In each case, this test was administered at age
11, so it would have been administered in 1946, 1958, and 1970. This is
an interesting quirk of the data that will be worth looking into in the
future.

\end{document}
